\documentclass[10pt,a4paper]{article}

% ------------------------------------------------------------------
%  NeetCode 150 -- Python revision notes
%  Compile with pdfLaTeX (works on Overleaf as-is).
%
%  HOW TO EDIT
%   * Each category lives in its own file under problems/.
%   * Every problem is one \begin{problem}...\end{problem} block:
%       \begin{problem}{<LeetCode no.>}{<Title>}{<Easy|Medium|Hard>}
%         \summary{...}  \keyidea{...}  \complexity{...}
%         \begin{code} ...python... \end{code}
%       \end{problem}
%   * Add your own notes with \note{...} inside a problem.
%   * The cheat sheet is in cheatsheet.tex.
%   * Colours, fonts and code style are all set in this preamble.
% ------------------------------------------------------------------

\usepackage[margin=16mm,top=15mm,bottom=16mm]{geometry}
\usepackage[T1]{fontenc}
\usepackage[scaled=0.92]{helvet}
\renewcommand{\familydefault}{\sfdefault}
\usepackage{textcomp}
\usepackage[dvipsnames]{xcolor}
\usepackage{listings}
\usepackage{multicol}
\usepackage{titlesec}
\usepackage{fancyhdr}
\usepackage{enumitem}
\usepackage{needspace}
\usepackage[hidelinks]{hyperref}
\usepackage{bookmark}

% ---------- colours ----------
\definecolor{ink}{HTML}{1D2330}
\definecolor{muted}{HTML}{5B6475}
\definecolor{accent}{HTML}{2F5D8A}
\definecolor{rule}{HTML}{D9DEE7}
\definecolor{codebg}{HTML}{F6F7F9}
\definecolor{easy}{HTML}{2E7D4F}
\definecolor{medium}{HTML}{B26A00}
\definecolor{hard}{HTML}{B3261E}
\definecolor{kw}{HTML}{7A3E9D}
\definecolor{bi}{HTML}{2F5D8A}
\definecolor{str}{HTML}{2E7D4F}
\definecolor{cmt}{HTML}{8A93A3}
\color{ink}
\setlength{\parindent}{0pt}

% ---------- code style ----------
\lstdefinestyle{py}{
  language=Python,
  basicstyle=\ttfamily\footnotesize,
  keywordstyle=\color{kw},
  keywordstyle=[2]\color{bi},
  morekeywords=[2]{len,range,enumerate,zip,min,max,sum,sorted,set,dict,list,tuple,
    int,str,bool,float,any,all,abs,iter,next,ord,map,divmod,reversed,
    List,Optional,defaultdict,Counter,deque,heapq,bisect,functools,self},
  morekeywords={None,True,False,nonlocal},
  stringstyle=\color{str},
  commentstyle=\color{cmt}\itshape,
  backgroundcolor=\color{codebg},
  frame=single, rulecolor=\color{rule}, framesep=5pt,
  xleftmargin=5pt, xrightmargin=5pt,
  columns=fullflexible, keepspaces=true, upquote=true,
  showstringspaces=false, breaklines=true,
  breakatwhitespace=true, tabsize=4,
  literate={~}{{\raisebox{-0.55ex}{\char126}}}1,
  aboveskip=4pt, belowskip=2pt
}
\lstnewenvironment{code}{\lstset{style=py}}{}
\lstnewenvironment{smallcode}{\lstset{style=py,basicstyle=\ttfamily\scriptsize,frame=none,xleftmargin=4pt,xrightmargin=4pt,framexleftmargin=4pt,framexrightmargin=4pt,framextopmargin=3pt,framexbottommargin=3pt}}{}

% ---------- headings ----------
\titleformat{\section}{\Large\bfseries\color{accent}}{}{0pt}{}[\vspace{-2pt}]
\titlespacing*{\section}{0pt}{0pt}{8pt}
\setcounter{secnumdepth}{0}
\setcounter{tocdepth}{2}

% ---------- header / footer ----------
\pagestyle{fancy}
\fancyhf{}
\renewcommand{\headrulewidth}{0pt}
\renewcommand{\footrulewidth}{0.4pt}
\renewcommand{\footrule}{{\color{rule}\hrule width\headwidth height\footrulewidth}}
\fancyfoot[L]{\scriptsize\color{muted}NeetCode 150 $\cdot$ Python revision}
\fancyfoot[R]{\scriptsize\color{muted}\leftmark\quad\thepage}
\renewcommand{\sectionmark}[1]{\markboth{#1}{}}

% ---------- difficulty tags ----------
\newcommand{\diff}[1]{%
  \def\tmp{#1}\def\e{Easy}\def\m{Medium}%
  \ifx\tmp\e\textcolor{easy}{\small\bfseries Easy}%
  \else\ifx\tmp\m\textcolor{medium}{\small\bfseries Medium}%
  \else\textcolor{hard}{\small\bfseries Hard}\fi\fi}

% short tag used in the contents
\newcommand{\tocdiff}[1]{%
  \def\tmp{#1}\def\e{Easy}\def\m{Medium}%
  \ifx\tmp\e\textcolor{easy}{\bfseries E}%
  \else\ifx\tmp\m\textcolor{medium}{\bfseries M}%
  \else\textcolor{hard}{\bfseries H}\fi\fi}

% ---------- problem block ----------
% \begin{problem}{number}{title}{difficulty}
\newenvironment{problem}[3]{%
  \Needspace{12\baselineskip}%
  \phantomsection
  \addcontentsline{toc}{subsection}{#1. #2 \texorpdfstring{\protect\tocdiff{#3}}{(#3)}}%
  \label{p#1}%
  \par\noindent{\large\bfseries #1. #2}\enspace\diff{#3}\par\vspace{2pt}%
  \small
}{\par\normalsize\vspace{10pt}}

\newcommand{\summary}[1]{\noindent\textbf{Problem:} #1\par}
\newcommand{\keyidea}[1]{\noindent\textbf{Key idea:} #1\par}
\newcommand{\complexity}[1]{\noindent{\color{muted}\textbf{Complexity:} #1}\par}
\newcommand{\note}[1]{\noindent{\color{accent}\textbf{My note:} #1}\par}

% ---------- cheat-sheet block ----------
% \cheat{category}{spot it}{core move}{traps}{template code (smallcode env)}
\newcommand{\cheatentry}[4]{%
  \textcolor{accent}{\bfseries #1}\par\vspace{2pt}
  \footnotesize
  \textbf{Spot it:} #2\par\vspace{1pt}
  \textbf{Core move:} #3\par\vspace{1pt}
  \textbf{Traps:} #4\par}

% ---------- contents formatting ----------
\makeatletter
\renewcommand{\l@section}[2]{\par\vspace{4pt}\noindent{\small\bfseries\color{accent}#1}\nobreak\par}
\renewcommand{\l@subsection}[2]{\noindent{\scriptsize #1}\par}
\makeatother

\hypersetup{
  pdftitle={NeetCode 150 -- Python Revision},
  pdfauthor={Claude Opus 5.5},
  bookmarksopen=false
}

\title{NeetCode 150}
\author{Claude Opus 5.5}
\date{}

\begin{document}

% ================= COVER =================
\thispagestyle{empty}
\pdfbookmark[0]{Cover \& how to use}{cover}
\vspace*{25mm}
{\fontsize{30}{36}\selectfont\bfseries NeetCode 150\par}
\vspace{4pt}
{\large\color{muted} Python revision notes $\cdot$ summary, key idea, complexity and solution per problem\par}
\vspace{6pt}
{\color{muted} Author: Claude Opus 5.5\par}
\vspace{10mm}
{\color{rule}\hrule}
\vspace{4pt}
\noindent
\begin{tabular}{@{}p{0.18\linewidth}p{0.18\linewidth}p{0.18\linewidth}p{0.18\linewidth}p{0.18\linewidth}@{}}
{\scriptsize\color{muted}Problems} & {\scriptsize\color{muted}Categories} & {\scriptsize\color{muted}Easy} & {\scriptsize\color{muted}Medium} & {\scriptsize\color{muted}Hard}\\
{\LARGE\bfseries 150} & {\LARGE\bfseries 18} & \textcolor{easy}{\LARGE\bfseries 28} & \textcolor{medium}{\LARGE\bfseries 100} & \textcolor{hard}{\LARGE\bfseries 22}
\end{tabular}

\vspace{10mm}
{\large\bfseries How to revise with this\par}
\begin{itemize}[leftmargin=12pt,itemsep=2pt]\small
  \item Read the \textbf{problem} line and try to recall the approach before reading the \textbf{key idea}. If it comes to you, skim the code; if not, read it slowly and re-derive why it works.
  \item The \textbf{pattern cheat sheet} is the fastest pass the night before a test: one block per category with the trigger, the template and the classic traps.
  \item The contents and PDF bookmarks link to every category and problem.
  \item Problem summaries are paraphrased, not copied from LeetCode; check the original statement for exact constraints.
  \item Add your own reminders inside any problem with \verb|\note{...}|.
\end{itemize}

\vspace{4pt}
{\large\bfseries Shared definitions assumed by the solutions\par}
{\small Solutions use LeetCode's signatures. These imports and node classes are provided by LeetCode and omitted from each problem.\par}
\begin{code}
from typing import List, Optional
from collections import defaultdict, Counter, deque
import heapq, bisect, functools

class ListNode:
    def __init__(self, val=0, next=None):
        self.val, self.next = val, next

class TreeNode:
    def __init__(self, val=0, left=None, right=None):
        self.val, self.left, self.right = val, left, right
\end{code}

\vspace{4pt}
{\footnotesize\color{muted} Problem list follows the classic NeetCode 150 grouping (18 categories); the live site may differ slightly. Every solution was run against test cases (LeetCode examples plus edge cases) before compiling; all passed.\par}

\newpage
% ================= CHEAT SHEET =================
\markboth{Cheat sheet}{}
\pdfbookmark[0]{Pattern cheat sheet}{cheat}
\section*{Pattern cheat sheet}
{\footnotesize\color{muted} One block per category: when to reach for it, the core move, a template, and the traps.\par}
\vspace{6pt}

% --- Arrays & Hashing ---
\noindent\begin{minipage}[t]{0.51\linewidth}
\cheatentry{Arrays \& Hashing}{Need O(1) lookup, counting, grouping, or 'have I seen this?'}{Hash set / hash map / Counter. Group by a canonical key (sorted string, count tuple). Prefix and suffix products/sums.}{Lists are unhashable: use tuple(count) as a key. Sequence-start trick: only count from n where n-1 is absent.}
\end{minipage}\hfill
\begin{minipage}[t]{0.46\linewidth}\vspace{-6pt}
\begin{smallcode}
seen = {}
for i, x in enumerate(nums):
    if target - x in seen: ...
    seen[x] = i
\end{smallcode}
\end{minipage}
\par\vspace{3pt}{\color{rule}\hrule}\vspace{6pt}

% --- Two Pointers ---
\noindent\begin{minipage}[t]{0.51\linewidth}
\cheatentry{Two Pointers}{Sorted input, pair/triplet sums, palindromes, or shrinking from both ends.}{l at start, r at end; move the pointer that can improve the answer.}{Skip duplicates after a match (3Sum). Move the shorter wall (Container, Trap).}
\end{minipage}\hfill
\begin{minipage}[t]{0.46\linewidth}\vspace{-6pt}
\begin{smallcode}
l, r = 0, len(a) - 1
while l < r:
    if cond: l += 1
    else: r -= 1
\end{smallcode}
\end{minipage}
\par\vspace{3pt}{\color{rule}\hrule}\vspace{6pt}

% --- Sliding Window ---
\noindent\begin{minipage}[t]{0.51\linewidth}
\cheatentry{Sliding Window}{Longest/shortest contiguous substring or subarray satisfying a condition.}{Grow r every step; shrink l while the window is invalid; record the answer.}{Fixed-size window: slide by adding r and removing r-k. Window max: monotonic deque.}
\end{minipage}\hfill
\begin{minipage}[t]{0.46\linewidth}\vspace{-6pt}
\begin{smallcode}
l = 0
for r in range(len(s)):
    add(s[r])
    while invalid(): remove(s[l]); l += 1
    best = max(best, r - l + 1)
\end{smallcode}
\end{minipage}
\par\vspace{3pt}{\color{rule}\hrule}\vspace{6pt}

% --- Stack ---
\noindent\begin{minipage}[t]{0.51\linewidth}
\cheatentry{Stack}{Matching pairs, 'next greater/warmer', expression evaluation, nested structure.}{Push openers / indices; pop when the current element resolves the top. Monotonic stack for next-greater.}{RPN division in Python: int(a / b), not a // b. Histogram: popped bar's start carries over.}
\end{minipage}\hfill
\begin{minipage}[t]{0.46\linewidth}\vspace{-6pt}
\begin{smallcode}
st = []
for i, x in enumerate(a):
    while st and a[st[-1]] < x:
        j = st.pop(); res[j] = i - j
    st.append(i)
\end{smallcode}
\end{minipage}
\par\vspace{3pt}{\color{rule}\hrule}\vspace{6pt}

% --- Binary Search ---
\noindent\begin{minipage}[t]{0.51\linewidth}
\cheatentry{Binary Search}{Sorted data, rotated arrays, or 'minimum X such that condition holds'.}{Halve the search space; on answer-space problems search over the answer value.}{Rotated array: one half is always sorted; check that half first. Median: partition the shorter array.}
\end{minipage}\hfill
\begin{minipage}[t]{0.46\linewidth}\vspace{-6pt}
\begin{smallcode}
l, r = lo, hi
while l <= r:
    m = (l + r) // 2
    if ok(m): res = m; r = m - 1
    else: l = m + 1
\end{smallcode}
\end{minipage}
\par\vspace{3pt}{\color{rule}\hrule}\vspace{6pt}

% --- Linked List ---
\noindent\begin{minipage}[t]{0.51\linewidth}
\cheatentry{Linked List}{In-place pointer rewiring, middle of list, cycles, k-th from end.}{Dummy head; slow/fast pointers; reverse with prev/cur/nxt.}{Find the Duplicate: array as a linked list (Floyd). LRU: hash map + doubly linked list with sentinels.}
\end{minipage}\hfill
\begin{minipage}[t]{0.46\linewidth}\vspace{-6pt}
\begin{smallcode}
prev, cur = None, head
while cur:
    nxt = cur.next; cur.next = prev
    prev, cur = cur, nxt
\end{smallcode}
\end{minipage}
\par\vspace{3pt}{\color{rule}\hrule}\vspace{6pt}

% --- Trees ---
\noindent\begin{minipage}[t]{0.51\linewidth}
\cheatentry{Trees}{Anything recursive over a hierarchy: depth, paths, validation, traversal order.}{DFS returning a value up (height, best arm), or BFS by level with a queue.}{BST: in-order is sorted; validate with (lo, hi) bounds, not just parent. Max path sum: clamp arms at 0.}
\end{minipage}\hfill
\begin{minipage}[t]{0.46\linewidth}\vspace{-6pt}
\begin{smallcode}
def dfs(node):
    if not node: return 0
    l, r = dfs(node.left), dfs(node.right)
    # update global answer with l, r
    return 1 + max(l, r)
\end{smallcode}
\end{minipage}
\par\vspace{3pt}{\color{rule}\hrule}\vspace{6pt}

% --- Tries ---
\noindent\begin{minipage}[t]{0.51\linewidth}
\cheatentry{Tries}{Prefix queries, word dictionaries, searching many words at once.}{Node = children dict + end flag. Walk character by character.}{Word Search II: prune found words and empty branches or it times out.}
\end{minipage}\hfill
\begin{minipage}[t]{0.46\linewidth}\vspace{-6pt}
\begin{smallcode}
node = root
for c in word:
    node = node.children.setdefault(c, TrieNode())
node.end = True
\end{smallcode}
\end{minipage}
\par\vspace{3pt}{\color{rule}\hrule}\vspace{6pt}

% --- Heap / Priority Queue ---
\noindent\begin{minipage}[t]{0.51\linewidth}
\cheatentry{Heap / Priority Queue}{k largest/smallest, running median, merging sorted streams, scheduling.}{heapq is a min-heap. For max-heap push negatives. Keep a heap of size k for top-k.}{Median: two heaps (max-heap low half, min-heap high half). Tie-breakers in tuples avoid comparing nodes.}
\end{minipage}\hfill
\begin{minipage}[t]{0.46\linewidth}\vspace{-6pt}
\begin{smallcode}
h = []
for x in nums:
    heapq.heappush(h, x)
    if len(h) > k: heapq.heappop(h)
# h[0] is k-th largest
\end{smallcode}
\end{minipage}
\par\vspace{3pt}{\color{rule}\hrule}\vspace{6pt}

% --- Backtracking ---
\noindent\begin{minipage}[t]{0.51\linewidth}
\cheatentry{Backtracking}{'Return all' combinations, subsets, permutations, placements.}{Choose, recurse, un-choose. Sort first when duplicates must be skipped.}{Append a copy (cur[:]). Skip equal neighbours at the same depth for duplicate inputs.}
\end{minipage}\hfill
\begin{minipage}[t]{0.46\linewidth}\vspace{-6pt}
\begin{smallcode}
def dfs(i):
    if done: res.append(cur[:]); return
    cur.append(x); dfs(i + 1); cur.pop()
    dfs(i + 1)
\end{smallcode}
\end{minipage}
\par\vspace{3pt}{\color{rule}\hrule}\vspace{6pt}

% --- Graphs ---
\noindent\begin{minipage}[t]{0.51\linewidth}
\cheatentry{Graphs}{Grids, connectivity, islands, dependencies, shortest steps in unweighted graphs.}{DFS/BFS flood fill; multi-source BFS for 'minutes/distance'; Kahn's topological sort; Union-Find.}{Mark visited when enqueuing, not when popping. Reverse the flow when it is easier (Pacific Atlantic, Surrounded Regions).}
\end{minipage}\hfill
\begin{minipage}[t]{0.46\linewidth}\vspace{-6pt}
\begin{smallcode}
q = deque(sources)
while q:
    for _ in range(len(q)):
        r, c = q.popleft()
        for nr, nc in neighbours(r, c): ...
    steps += 1
\end{smallcode}
\end{minipage}
\par\vspace{3pt}{\color{rule}\hrule}\vspace{6pt}

% --- Advanced Graphs ---
\noindent\begin{minipage}[t]{0.51\linewidth}
\cheatentry{Advanced Graphs}{Weighted shortest paths, minimum spanning tree, ordering from constraints, Euler paths.}{Dijkstra (heap), Prim (heap), Bellman-Ford (k rounds), Kahn (topo), Hierholzer (Euler path).}{K stops: Bellman-Ford with a copy each round. Alien Dictionary: a word before its own prefix is invalid.}
\end{minipage}\hfill
\begin{minipage}[t]{0.46\linewidth}\vspace{-6pt}
\begin{smallcode}
h = [(0, src)]; dist = {}
while h:
    d, u = heapq.heappop(h)
    if u in dist: continue
    dist[u] = d
    for v, w in adj[u]: heapq.heappush(h, (d + w, v))
\end{smallcode}
\end{minipage}
\par\vspace{3pt}{\color{rule}\hrule}\vspace{6pt}

% --- 1-D Dynamic Programming ---
\noindent\begin{minipage}[t]{0.51\linewidth}
\cheatentry{1-D Dynamic Programming}{Count ways / min cost / max value where the answer at i depends on a few earlier states.}{Define dp[i], write the recurrence, then compress to two variables if possible.}{Palindromes: expand around centres. Max product: track min and max. LIS: patience sorting + bisect.}
\end{minipage}\hfill
\begin{minipage}[t]{0.46\linewidth}\vspace{-6pt}
\begin{smallcode}
a, b = base0, base1
for x in arr:
    a, b = b, max(b, a + x)   # House Robber
\end{smallcode}
\end{minipage}
\par\vspace{3pt}{\color{rule}\hrule}\vspace{6pt}

% --- 2-D Dynamic Programming ---
\noindent\begin{minipage}[t]{0.51\linewidth}
\cheatentry{2-D Dynamic Programming}{Two strings or a grid; choices at each (i, j); knapsack-style counting.}{dp[i][j] over suffixes (or prefixes); fill bottom-up or memoise top-down.}{Coin Change II: coins in the outer loop (combinations, not permutations). Burst Balloons: pick the LAST balloon.}
\end{minipage}\hfill
\begin{minipage}[t]{0.46\linewidth}\vspace{-6pt}
\begin{smallcode}
for i in range(m - 1, -1, -1):
    for j in range(n - 1, -1, -1):
        dp[i][j] = f(dp[i+1][j],
                     dp[i][j+1],
                     dp[i+1][j+1])
\end{smallcode}
\end{minipage}
\par\vspace{3pt}{\color{rule}\hrule}\vspace{6pt}

% --- Greedy ---
\noindent\begin{minipage}[t]{0.51\linewidth}
\cheatentry{Greedy}{A locally best move is provably safe: max subarray, reachability, partitioning.}{Keep one running quantity (sum, goal, reach, open-count range) and update it each step.}{Gas Station: total check first, then reset start. Parenthesis with *: track a [lo, hi] range.}
\end{minipage}\hfill
\begin{minipage}[t]{0.46\linewidth}\vspace{-6pt}
\begin{smallcode}
cur = best = nums[0]
for x in nums[1:]:
    cur = max(x, cur + x)
    best = max(best, cur)
\end{smallcode}
\end{minipage}
\par\vspace{3pt}{\color{rule}\hrule}\vspace{6pt}

% --- Intervals ---
\noindent\begin{minipage}[t]{0.51\linewidth}
\cheatentry{Intervals}{Overlaps, merging, inserting, rooms needed, queries over ranges.}{Sort by start (or end); compare each interval's start with the previous end.}{Non-overlapping: on conflict keep the smaller end. Meeting Rooms II: separate sorted starts and ends.}
\end{minipage}\hfill
\begin{minipage}[t]{0.46\linewidth}\vspace{-6pt}
\begin{smallcode}
intervals.sort()
res = [intervals[0]]
for s, e in intervals[1:]:
    if s <= res[-1][1]:
        res[-1][1] = max(res[-1][1], e)
    else:
        res.append([s, e])
\end{smallcode}
\end{minipage}
\par\vspace{3pt}{\color{rule}\hrule}\vspace{6pt}

% --- Math & Geometry ---
\noindent\begin{minipage}[t]{0.51\linewidth}
\cheatentry{Math \& Geometry}{Matrix manipulation, digit arithmetic, cycles in number sequences.}{Transpose + reverse to rotate; boundary pointers for spirals; fast power by squaring.}{Set Matrix Zeroes in O(1) space: first row/col as markers plus one flag. Negative exponents: invert x.}
\end{minipage}\hfill
\begin{minipage}[t]{0.46\linewidth}\vspace{-6pt}
\begin{smallcode}
top, bottom, left, right = 0, R - 1, 0, C - 1
while top <= bottom and left <= right: ...
\end{smallcode}
\end{minipage}
\par\vspace{3pt}{\color{rule}\hrule}\vspace{6pt}

% --- Bit Manipulation ---
\noindent\begin{minipage}[t]{0.51\linewidth}
\cheatentry{Bit Manipulation}{XOR cancellation, counting bits, arithmetic without + or -.}{x \textasciicircum{} x = 0; n \& (n - 1) clears the lowest set bit; shifts move bits.}{Python ints are unbounded: mask with 0xFFFFFFFF and convert back for negatives (Sum of Two Integers).}
\end{minipage}\hfill
\begin{minipage}[t]{0.46\linewidth}\vspace{-6pt}
\begin{smallcode}
res = 0
for x in nums: res ^= x
\end{smallcode}
\end{minipage}
\par\vspace{3pt}{\color{rule}\hrule}\vspace{6pt}


\newpage
% ================= CONTENTS =================
\markboth{Contents}{}
\pdfbookmark[0]{Contents}{contents}
\section*{Contents}
\begin{multicols}{2}
\makeatletter\@starttoc{toc}\makeatother
\end{multicols}

% ================= PROBLEMS =================
\newpage
% ===== Arrays & Hashing (9 problems) =====
\section{Arrays \& Hashing \texorpdfstring{{\normalsize\color{muted}$\cdot$ 9 problems}}{- 9 problems}}
\markboth{Arrays \& Hashing}{}

\begin{problem}{217}{Contains Duplicate}{Easy}
\summary{Return True if any value appears at least twice in the array.}
\keyidea{A set remembers what you've seen; a repeat hit means a duplicate.}
\complexity{O(n) time, O(n) space}
\begin{code}
class Solution:
    def containsDuplicate(self, nums: List[int]) -> bool:
        seen = set()
        for n in nums:
            if n in seen:
                return True
            seen.add(n)
        return False
\end{code}
\end{problem}

\begin{problem}{242}{Valid Anagram}{Easy}
\summary{Check whether string t is a rearrangement of string s.}
\keyidea{Same length and identical character counts.}
\complexity{O(n) time, O(1) space (26 letters)}
\begin{code}
class Solution:
    def isAnagram(self, s: str, t: str) -> bool:
        if len(s) != len(t):
            return False
        count = {}
        for a, b in zip(s, t):
            count[a] = count.get(a, 0) + 1
            count[b] = count.get(b, 0) - 1
        return all(v == 0 for v in count.values())
\end{code}
\end{problem}

\begin{problem}{1}{Two Sum}{Easy}
\summary{Return indices of the two numbers that add up to target (exactly one answer exists).}
\keyidea{For each number, look up its complement (target - n) in a value-\textgreater{}index map built so far.}
\complexity{O(n) time, O(n) space}
\begin{code}
class Solution:
    def twoSum(self, nums: List[int], target: int) -> List[int]:
        idx = {}
        for i, n in enumerate(nums):
            if target - n in idx:
                return [idx[target - n], i]
            idx[n] = i
\end{code}
\end{problem}

\begin{problem}{49}{Group Anagrams}{Medium}
\summary{Group the given strings so that anagrams sit in the same group.}
\keyidea{Key each word by its 26-letter count tuple (or sorted string); bucket into a dict.}
\complexity{O(n * k) time, O(n * k) space}
\begin{code}
class Solution:
    def groupAnagrams(self, strs: List[str]) -> List[List[str]]:
        groups = defaultdict(list)
        for s in strs:
            count = [0] * 26
            for c in s:
                count[ord(c) - ord('a')] += 1
            groups[tuple(count)].append(s)
        return list(groups.values())
\end{code}
\end{problem}

\begin{problem}{347}{Top K Frequent Elements}{Medium}
\summary{Return the k values that occur most often.}
\keyidea{Bucket sort: index = frequency, then walk buckets from high to low.}
\complexity{O(n) time, O(n) space}
\begin{code}
class Solution:
    def topKFrequent(self, nums: List[int], k: int) -> List[int]:
        count = Counter(nums)
        buckets = [[] for _ in range(len(nums) + 1)]
        for n, c in count.items():
            buckets[c].append(n)
        res = []
        for i in range(len(buckets) - 1, 0, -1):
            for n in buckets[i]:
                res.append(n)
                if len(res) == k:
                    return res
\end{code}
\end{problem}

\begin{problem}{238}{Product of Array Except Self}{Medium}
\summary{For each index, output the product of every other element, without division, in O(n).}
\keyidea{Answer = prefix product (left of i) x suffix product (right of i); do it in two passes.}
\complexity{O(n) time, O(1) extra space (output excluded)}
\begin{code}
class Solution:
    def productExceptSelf(self, nums: List[int]) -> List[int]:
        res = [1] * len(nums)
        prefix = 1
        for i in range(len(nums)):
            res[i] = prefix
            prefix *= nums[i]
        suffix = 1
        for i in range(len(nums) - 1, -1, -1):
            res[i] *= suffix
            suffix *= nums[i]
        return res
\end{code}
\end{problem}

\begin{problem}{36}{Valid Sudoku}{Medium}
\summary{Decide if a partially filled 9x9 board breaks no Sudoku rule (rows, columns, 3x3 boxes).}
\keyidea{Three sets of sets: rows[r], cols[c], boxes[(r//3, c//3)]. Any repeat is invalid.}
\complexity{O(81) = O(1) time and space}
\begin{code}
class Solution:
    def isValidSudoku(self, board: List[List[str]]) -> bool:
        rows, cols, boxes = defaultdict(set), defaultdict(set), defaultdict(set)
        for r in range(9):
            for c in range(9):
                v = board[r][c]
                if v == ".":
                    continue
                if v in rows[r] or v in cols[c] or v in boxes[(r // 3, c // 3)]:
                    return False
                rows[r].add(v)
                cols[c].add(v)
                boxes[(r // 3, c // 3)].add(v)
        return True
\end{code}
\end{problem}

\begin{problem}{271}{Encode and Decode Strings}{Medium}
\summary{Design encode(list of strings) -\textgreater{} one string and decode back, for any characters.}
\keyidea{Length-prefix each word: '4\#leet'. Read the number up to '\#', then slice that many chars.}
\complexity{O(total chars) time and space}
\begin{code}
class Solution:
    def encode(self, strs: List[str]) -> str:
        return "".join(f"{len(s)}#{s}" for s in strs)

    def decode(self, s: str) -> List[str]:
        res, i = [], 0
        while i < len(s):
            j = i
            while s[j] != "#":
                j += 1
            length = int(s[i:j])
            res.append(s[j + 1 : j + 1 + length])
            i = j + 1 + length
        return res
\end{code}
\end{problem}

\begin{problem}{128}{Longest Consecutive Sequence}{Medium}
\summary{Length of the longest run of consecutive integers in an unsorted array, in O(n).}
\keyidea{Put all in a set; only start counting from n where n-1 is absent (a sequence start).}
\complexity{O(n) time, O(n) space}
\begin{code}
class Solution:
    def longestConsecutive(self, nums: List[int]) -> int:
        s = set(nums)
        best = 0
        for n in s:
            if n - 1 not in s:
                length = 1
                while n + length in s:
                    length += 1
                best = max(best, length)
        return best
\end{code}
\end{problem}

\newpage
% ===== Two Pointers (5 problems) =====
\section{Two Pointers {\normalsize\color{muted}$\cdot$ 5 problems}}
\markboth{Two Pointers}{}

\begin{problem}{125}{Valid Palindrome}{Easy}
\summary{After keeping only alphanumerics and lower-casing, is the string a palindrome?}
\keyidea{Pointers from both ends, skip non-alphanumerics, compare lowercase.}
\complexity{O(n) time, O(1) space}
\begin{code}
class Solution:
    def isPalindrome(self, s: str) -> bool:
        l, r = 0, len(s) - 1
        while l < r:
            while l < r and not s[l].isalnum():
                l += 1
            while l < r and not s[r].isalnum():
                r -= 1
            if s[l].lower() != s[r].lower():
                return False
            l, r = l + 1, r - 1
        return True
\end{code}
\end{problem}

\begin{problem}{167}{Two Sum II - Input Array Is Sorted}{Medium}
\summary{Sorted array; return 1-indexed positions of two numbers summing to target, O(1) space.}
\keyidea{l at start, r at end: sum too big -\textgreater{} r--, too small -\textgreater{} l++.}
\complexity{O(n) time, O(1) space}
\begin{code}
class Solution:
    def twoSum(self, numbers: List[int], target: int) -> List[int]:
        l, r = 0, len(numbers) - 1
        while l < r:
            s = numbers[l] + numbers[r]
            if s > target:
                r -= 1
            elif s < target:
                l += 1
            else:
                return [l + 1, r + 1]
\end{code}
\end{problem}

\begin{problem}{15}{3Sum}{Medium}
\summary{Find all unique triplets that sum to zero.}
\keyidea{Sort; fix a, then two-pointer for the pair. Skip duplicate a's and duplicate l's.}
\complexity{O(n\textasciicircum{}2) time, O(1) extra space (besides sort)}
\begin{code}
class Solution:
    def threeSum(self, nums: List[int]) -> List[List[int]]:
        nums.sort()
        res = []
        for i, a in enumerate(nums):
            if a > 0:
                break
            if i > 0 and a == nums[i - 1]:
                continue
            l, r = i + 1, len(nums) - 1
            while l < r:
                s = a + nums[l] + nums[r]
                if s > 0:
                    r -= 1
                elif s < 0:
                    l += 1
                else:
                    res.append([a, nums[l], nums[r]])
                    l += 1
                    while l < r and nums[l] == nums[l - 1]:
                        l += 1
        return res
\end{code}
\end{problem}

\begin{problem}{11}{Container With Most Water}{Medium}
\summary{Pick two vertical lines that, with the x-axis, hold the most water.}
\keyidea{Start widest; always move the shorter wall inward (the taller one can't help).}
\complexity{O(n) time, O(1) space}
\begin{code}
class Solution:
    def maxArea(self, height: List[int]) -> int:
        l, r, best = 0, len(height) - 1, 0
        while l < r:
            best = max(best, (r - l) * min(height[l], height[r]))
            if height[l] < height[r]:
                l += 1
            else:
                r -= 1
        return best
\end{code}
\end{problem}

\begin{problem}{42}{Trapping Rain Water}{Hard}
\summary{Given bar heights, compute how much rain water is trapped between them.}
\keyidea{Water at i = min(maxLeft, maxRight) - h[i]. Two pointers: process the side with the smaller max.}
\complexity{O(n) time, O(1) space}
\begin{code}
class Solution:
    def trap(self, height: List[int]) -> int:
        if not height:
            return 0
        l, r = 0, len(height) - 1
        leftMax, rightMax = height[l], height[r]
        res = 0
        while l < r:
            if leftMax < rightMax:
                l += 1
                leftMax = max(leftMax, height[l])
                res += leftMax - height[l]
            else:
                r -= 1
                rightMax = max(rightMax, height[r])
                res += rightMax - height[r]
        return res
\end{code}
\end{problem}

\newpage
% ===== Sliding Window (6 problems) =====
\section{Sliding Window \texorpdfstring{{\normalsize\color{muted}$\cdot$ 6 problems}}{- 6 problems}}
\markboth{Sliding Window}{}

\begin{problem}{121}{Best Time to Buy and Sell Stock}{Easy}
\summary{One buy then one later sell; maximise profit (0 if none).}
\keyidea{Track the minimum price so far; profit = today - min.}
\complexity{O(n) time, O(1) space}
\begin{code}
class Solution:
    def maxProfit(self, prices: List[int]) -> int:
        lo, best = float("inf"), 0
        for p in prices:
            lo = min(lo, p)
            best = max(best, p - lo)
        return best
\end{code}
\end{problem}

\begin{problem}{3}{Longest Substring Without Repeating Characters}{Medium}
\summary{Length of the longest substring with all distinct characters.}
\keyidea{Window with a set; on a repeat, shrink from the left until the char is gone.}
\complexity{O(n) time, O(alphabet) space}
\begin{code}
class Solution:
    def lengthOfLongestSubstring(self, s: str) -> int:
        seen = set()
        l = best = 0
        for r in range(len(s)):
            while s[r] in seen:
                seen.remove(s[l])
                l += 1
            seen.add(s[r])
            best = max(best, r - l + 1)
        return best
\end{code}
\end{problem}

\begin{problem}{424}{Longest Repeating Character Replacement}{Medium}
\summary{You may change up to k characters; longest substring that can become all one letter.}
\keyidea{Window valid while (length - count of most frequent char) \textless{}= k. maxf never needs decreasing.}
\complexity{O(n) time, O(26) space}
\begin{code}
class Solution:
    def characterReplacement(self, s: str, k: int) -> int:
        count = {}
        l = maxf = best = 0
        for r in range(len(s)):
            count[s[r]] = count.get(s[r], 0) + 1
            maxf = max(maxf, count[s[r]])
            while (r - l + 1) - maxf > k:
                count[s[l]] -= 1
                l += 1
            best = max(best, r - l + 1)
        return best
\end{code}
\end{problem}

\begin{problem}{567}{Permutation in String}{Medium}
\summary{Does s2 contain some rearrangement of s1 as a contiguous substring?}
\keyidea{Fixed-size window of len(s1); compare 26-letter counts as the window slides.}
\complexity{O(n) time, O(26) space}
\begin{code}
class Solution:
    def checkInclusion(self, s1: str, s2: str) -> bool:
        if len(s1) > len(s2):
            return False
        need, have = [0] * 26, [0] * 26
        for i in range(len(s1)):
            need[ord(s1[i]) - 97] += 1
            have[ord(s2[i]) - 97] += 1
        if need == have:
            return True
        for r in range(len(s1), len(s2)):
            have[ord(s2[r]) - 97] += 1
            have[ord(s2[r - len(s1)]) - 97] -= 1
            if need == have:
                return True
        return False
\end{code}
\end{problem}

\begin{problem}{76}{Minimum Window Substring}{Hard}
\summary{Smallest substring of s containing every character of t (with multiplicity).}
\keyidea{Expand r until all needs met (have == need), then shrink l while still valid, recording best.}
\complexity{O(\textbar{}s\textbar{} + \textbar{}t\textbar{}) time, O(alphabet) space}
\begin{code}
class Solution:
    def minWindow(self, s: str, t: str) -> str:
        if not t:
            return ""
        countT, window = Counter(t), {}
        have, need = 0, len(countT)
        res, resLen = (-1, -1), float("inf")
        l = 0
        for r, c in enumerate(s):
            window[c] = window.get(c, 0) + 1
            if c in countT and window[c] == countT[c]:
                have += 1
            while have == need:
                if r - l + 1 < resLen:
                    res, resLen = (l, r), r - l + 1
                window[s[l]] -= 1
                if s[l] in countT and window[s[l]] < countT[s[l]]:
                    have -= 1
                l += 1
        l, r = res
        return s[l : r + 1] if resLen != float("inf") else ""
\end{code}
\end{problem}

\begin{problem}{239}{Sliding Window Maximum}{Hard}
\summary{Return the maximum of every window of size k.}
\keyidea{Monotonic decreasing deque of indices; front is the max; pop expired front.}
\complexity{O(n) time, O(k) space}
\begin{code}
class Solution:
    def maxSlidingWindow(self, nums: List[int], k: int) -> List[int]:
        q, res = deque(), []
        for r, n in enumerate(nums):
            while q and nums[q[-1]] < n:
                q.pop()
            q.append(r)
            if q[0] <= r - k:
                q.popleft()
            if r >= k - 1:
                res.append(nums[q[0]])
        return res
\end{code}
\end{problem}

\newpage
% ===== Stack (7 problems) =====
\section{Stack \texorpdfstring{{\normalsize\color{muted}$\cdot$ 7 problems}}{- 7 problems}}
\markboth{Stack}{}

\begin{problem}{20}{Valid Parentheses}{Easy}
\summary{Are brackets ()[]\{\} correctly matched and nested?}
\keyidea{Push openers; on a closer, the top must be its matching opener.}
\complexity{O(n) time, O(n) space}
\begin{code}
class Solution:
    def isValid(self, s: str) -> bool:
        pairs = {")": "(", "]": "[", "}": "{"}
        stack = []
        for c in s:
            if c in pairs:
                if not stack or stack[-1] != pairs[c]:
                    return False
                stack.pop()
            else:
                stack.append(c)
        return not stack
\end{code}
\end{problem}

\begin{problem}{155}{Min Stack}{Medium}
\summary{Stack supporting push, pop, top and getMin, all O(1).}
\keyidea{Second stack storing the min at each depth.}
\complexity{O(1) per op, O(n) space}
\begin{code}
class MinStack:
    def __init__(self):
        self.stack, self.mins = [], []

    def push(self, val: int) -> None:
        self.stack.append(val)
        self.mins.append(min(val, self.mins[-1] if self.mins else val))

    def pop(self) -> None:
        self.stack.pop()
        self.mins.pop()

    def top(self) -> int:
        return self.stack[-1]

    def getMin(self) -> int:
        return self.mins[-1]
\end{code}
\end{problem}

\begin{problem}{150}{Evaluate Reverse Polish Notation}{Medium}
\summary{Evaluate an expression in postfix notation (+ - * /, division truncates toward zero).}
\keyidea{Push numbers; on an operator pop b then a, push a op b. Use int(a / b) for truncation.}
\complexity{O(n) time, O(n) space}
\begin{code}
class Solution:
    def evalRPN(self, tokens: List[str]) -> int:
        st = []
        for t in tokens:
            if t in "+-*/":
                b, a = st.pop(), st.pop()
                if t == "+": st.append(a + b)
                elif t == "-": st.append(a - b)
                elif t == "*": st.append(a * b)
                else: st.append(int(a / b))
            else:
                st.append(int(t))
        return st[0]
\end{code}
\end{problem}

\begin{problem}{22}{Generate Parentheses}{Medium}
\summary{List every well-formed string of n pairs of parentheses.}
\keyidea{Backtrack: add '(' if open \textless{} n; add ')' if close \textless{} open.}
\complexity{O(4\textasciicircum{}n / sqrt(n)) time (Catalan), O(n) recursion}
\begin{code}
class Solution:
    def generateParenthesis(self, n: int) -> List[str]:
        res, cur = [], []
        def bt(op, cl):
            if op == cl == n:
                res.append("".join(cur))
                return
            if op < n:
                cur.append("("); bt(op + 1, cl); cur.pop()
            if cl < op:
                cur.append(")"); bt(op, cl + 1); cur.pop()
        bt(0, 0)
        return res
\end{code}
\end{problem}

\begin{problem}{739}{Daily Temperatures}{Medium}
\summary{For each day, how many days until a warmer temperature (0 if never)?}
\keyidea{Monotonic decreasing stack of indices; a warmer day resolves everything colder on the stack.}
\complexity{O(n) time, O(n) space}
\begin{code}
class Solution:
    def dailyTemperatures(self, temperatures: List[int]) -> List[int]:
        res = [0] * len(temperatures)
        st = []  # indices
        for i, t in enumerate(temperatures):
            while st and temperatures[st[-1]] < t:
                j = st.pop()
                res[j] = i - j
            st.append(i)
        return res
\end{code}
\end{problem}

\begin{problem}{853}{Car Fleet}{Medium}
\summary{Cars drive to target; a faster car that catches a slower one joins it. Count arriving fleets.}
\keyidea{Sort by position descending; compute arrival time. A car forms a new fleet only if it takes longer than the fleet ahead.}
\complexity{O(n log n) time, O(n) space}
\begin{code}
class Solution:
    def carFleet(self, target: int, position: List[int], speed: List[int]) -> int:
        cars = sorted(zip(position, speed), reverse=True)
        st = []
        for p, s in cars:
            t = (target - p) / s
            if not st or t > st[-1]:
                st.append(t)
        return len(st)
\end{code}
\end{problem}

\begin{problem}{84}{Largest Rectangle in Histogram}{Hard}
\summary{Area of the largest rectangle that fits under the histogram bars.}
\keyidea{Increasing stack of (start, height). When popping a taller bar, its rectangle ends here; the new bar inherits the popped start.}
\complexity{O(n) time, O(n) space}
\begin{code}
class Solution:
    def largestRectangleArea(self, heights: List[int]) -> int:
        best, st = 0, []  # (start index, height)
        for i, h in enumerate(heights):
            start = i
            while st and st[-1][1] > h:
                idx, height = st.pop()
                best = max(best, height * (i - idx))
                start = idx
            st.append((start, h))
        for idx, height in st:
            best = max(best, height * (len(heights) - idx))
        return best
\end{code}
\end{problem}

\newpage
% ===== Binary Search (7 problems) =====
\section{Binary Search \texorpdfstring{{\normalsize\color{muted}$\cdot$ 7 problems}}{- 7 problems}}
\markboth{Binary Search}{}

\begin{problem}{704}{Binary Search}{Easy}
\summary{Return the index of target in a sorted array, or -1.}
\keyidea{Classic l \textless{}= r loop with mid = l + (r - l) // 2.}
\complexity{O(log n) time, O(1) space}
\begin{code}
class Solution:
    def search(self, nums: List[int], target: int) -> int:
        l, r = 0, len(nums) - 1
        while l <= r:
            m = (l + r) // 2
            if nums[m] == target:
                return m
            if nums[m] < target:
                l = m + 1
            else:
                r = m - 1
        return -1
\end{code}
\end{problem}

\begin{problem}{74}{Search a 2D Matrix}{Medium}
\summary{Rows sorted and each row starts after the previous ends. Is target present?}
\keyidea{Treat it as one sorted array of m*n: index -\textgreater{} (i // n, i \% n).}
\complexity{O(log(m*n)) time, O(1) space}
\begin{code}
class Solution:
    def searchMatrix(self, matrix: List[List[int]], target: int) -> bool:
        m, n = len(matrix), len(matrix[0])
        l, r = 0, m * n - 1
        while l <= r:
            mid = (l + r) // 2
            v = matrix[mid // n][mid % n]
            if v == target:
                return True
            if v < target:
                l = mid + 1
            else:
                r = mid - 1
        return False
\end{code}
\end{problem}

\begin{problem}{875}{Koko Eating Bananas}{Medium}
\summary{Minimum eating speed k so all piles are finished within h hours (one pile per hour max).}
\keyidea{Binary search on the answer k in [1, max(piles)]; hours(k) = sum(ceil(p / k)).}
\complexity{O(n log max) time, O(1) space}
\begin{code}
class Solution:
    def minEatingSpeed(self, piles: List[int], h: int) -> int:
        l, r = 1, max(piles)
        res = r
        while l <= r:
            k = (l + r) // 2
            hours = sum((p + k - 1) // k for p in piles)
            if hours <= h:
                res = k
                r = k - 1
            else:
                l = k + 1
        return res
\end{code}
\end{problem}

\begin{problem}{153}{Find Minimum in Rotated Sorted Array}{Medium}
\summary{A sorted array of distinct values was rotated; find its minimum in O(log n).}
\keyidea{Compare mid with the right end: nums[m] \textgreater{} nums[r] means the min is right of m.}
\complexity{O(log n) time, O(1) space}
\begin{code}
class Solution:
    def findMin(self, nums: List[int]) -> int:
        l, r = 0, len(nums) - 1
        while l < r:
            m = (l + r) // 2
            if nums[m] > nums[r]:
                l = m + 1
            else:
                r = m
        return nums[l]
\end{code}
\end{problem}

\begin{problem}{33}{Search in Rotated Sorted Array}{Medium}
\summary{Find target's index in a rotated sorted array of distinct values, O(log n).}
\keyidea{One half is always sorted. Check if target lies inside the sorted half; else go to the other half.}
\complexity{O(log n) time, O(1) space}
\begin{code}
class Solution:
    def search(self, nums: List[int], target: int) -> int:
        l, r = 0, len(nums) - 1
        while l <= r:
            m = (l + r) // 2
            if nums[m] == target:
                return m
            if nums[l] <= nums[m]:            # left half sorted
                if nums[l] <= target < nums[m]:
                    r = m - 1
                else:
                    l = m + 1
            else:                              # right half sorted
                if nums[m] < target <= nums[r]:
                    l = m + 1
                else:
                    r = m - 1
        return -1
\end{code}
\end{problem}

\begin{problem}{981}{Time Based Key-Value Store}{Medium}
\summary{set(key, value, timestamp); get(key, t) returns the value with the largest timestamp \textless{}= t.}
\keyidea{Timestamps arrive increasing, so each key's list is sorted: binary search for the last ts \textless{}= t.}
\complexity{set O(1), get O(log n); O(n) space}
\begin{code}
class TimeMap:
    def __init__(self):
        self.store = defaultdict(list)  # key -> [(ts, value)]

    def set(self, key: str, value: str, timestamp: int) -> None:
        self.store[key].append((timestamp, value))

    def get(self, key: str, timestamp: int) -> str:
        arr, res = self.store[key], ""
        l, r = 0, len(arr) - 1
        while l <= r:
            m = (l + r) // 2
            if arr[m][0] <= timestamp:
                res = arr[m][1]
                l = m + 1
            else:
                r = m - 1
        return res
\end{code}
\end{problem}

\begin{problem}{4}{Median of Two Sorted Arrays}{Hard}
\summary{Median of two sorted arrays combined, in O(log(min(m, n))).}
\keyidea{Binary search a cut in the shorter array so left halves together hold half the elements and Aleft \textless{}= Bright, Bleft \textless{}= Aright.}
\complexity{O(log min(m, n)) time, O(1) space}
\begin{code}
class Solution:
    def findMedianSortedArrays(self, nums1: List[int], nums2: List[int]) -> float:
        A, B = nums1, nums2
        if len(A) > len(B):
            A, B = B, A
        total = len(A) + len(B)
        half = total // 2
        l, r = 0, len(A)
        while True:
            i = (l + r) // 2          # elements taken from A
            j = half - i              # elements taken from B
            Aleft = A[i - 1] if i > 0 else float("-inf")
            Aright = A[i] if i < len(A) else float("inf")
            Bleft = B[j - 1] if j > 0 else float("-inf")
            Bright = B[j] if j < len(B) else float("inf")
            if Aleft <= Bright and Bleft <= Aright:
                if total % 2:
                    return min(Aright, Bright)
                return (max(Aleft, Bleft) + min(Aright, Bright)) / 2
            if Aleft > Bright:
                r = i - 1
            else:
                l = i + 1
\end{code}
\end{problem}

\newpage
% ===== Linked List (11 problems) =====
\section{Linked List {\normalsize\color{muted}$\cdot$ 11 problems}}
\markboth{Linked List}{}

\begin{problem}{206}{Reverse Linked List}{Easy}
\summary{Reverse a singly linked list.}
\keyidea{Walk once, pointing each node back to prev.}
\complexity{O(n) time, O(1) space}
\begin{code}
class Solution:
    def reverseList(self, head: Optional[ListNode]) -> Optional[ListNode]:
        prev, cur = None, head
        while cur:
            nxt = cur.next
            cur.next = prev
            prev, cur = cur, nxt
        return prev
\end{code}
\end{problem}

\begin{problem}{21}{Merge Two Sorted Lists}{Easy}
\summary{Splice two sorted lists into one sorted list.}
\keyidea{Dummy head; attach the smaller node each step; append the leftover.}
\complexity{O(m + n) time, O(1) space}
\begin{code}
class Solution:
    def mergeTwoLists(self, list1: Optional[ListNode], list2: Optional[ListNode]) -> Optional[ListNode]:
        dummy = tail = ListNode()
        while list1 and list2:
            if list1.val <= list2.val:
                tail.next, list1 = list1, list1.next
            else:
                tail.next, list2 = list2, list2.next
            tail = tail.next
        tail.next = list1 or list2
        return dummy.next
\end{code}
\end{problem}

\begin{problem}{143}{Reorder List}{Medium}
\summary{Rearrange L0,L1,...,Ln into L0,Ln,L1,Ln-1,... in place.}
\keyidea{Three steps: find middle (slow/fast), reverse the second half, merge alternately.}
\complexity{O(n) time, O(1) space}
\begin{code}
class Solution:
    def reorderList(self, head: Optional[ListNode]) -> None:
        slow, fast = head, head.next
        while fast and fast.next:
            slow, fast = slow.next, fast.next.next
        second = slow.next
        slow.next = None
        prev = None
        while second:                      # reverse second half
            nxt = second.next
            second.next = prev
            prev, second = second, nxt
        first, second = head, prev
        while second:                      # merge
            t1, t2 = first.next, second.next
            first.next = second
            second.next = t1
            first, second = t1, t2
\end{code}
\end{problem}

\begin{problem}{19}{Remove Nth Node From End of List}{Medium}
\summary{Delete the n-th node from the end in one pass.}
\keyidea{Dummy head; move right pointer n ahead, then move both until right is None; left.next is the target.}
\complexity{O(n) time, O(1) space}
\begin{code}
class Solution:
    def removeNthFromEnd(self, head: Optional[ListNode], n: int) -> Optional[ListNode]:
        dummy = ListNode(0, head)
        left, right = dummy, head
        for _ in range(n):
            right = right.next
        while right:
            left, right = left.next, right.next
        left.next = left.next.next
        return dummy.next
\end{code}
\end{problem}

\begin{problem}{138}{Copy List with Random Pointer}{Medium}
\summary{Deep-copy a list where each node also has a random pointer.}
\keyidea{Pass 1: map old -\textgreater{} new node. Pass 2: wire next and random through the map.}
\complexity{O(n) time, O(n) space}
\begin{code}
class Solution:
    def copyRandomList(self, head: "Optional[Node]") -> "Optional[Node]":
        old_to_new = {None: None}
        cur = head
        while cur:
            old_to_new[cur] = Node(cur.val)
            cur = cur.next
        cur = head
        while cur:
            copy = old_to_new[cur]
            copy.next = old_to_new[cur.next]
            copy.random = old_to_new[cur.random]
            cur = cur.next
        return old_to_new[head]
\end{code}
\end{problem}

\begin{problem}{2}{Add Two Numbers}{Medium}
\summary{Two numbers stored as reversed digit lists; return their sum as a list.}
\keyidea{Add digit by digit with a carry; keep going while either list or carry remains.}
\complexity{O(max(m, n)) time, O(1) extra space}
\begin{code}
class Solution:
    def addTwoNumbers(self, l1: Optional[ListNode], l2: Optional[ListNode]) -> Optional[ListNode]:
        dummy = cur = ListNode()
        carry = 0
        while l1 or l2 or carry:
            v = (l1.val if l1 else 0) + (l2.val if l2 else 0) + carry
            carry, v = divmod(v, 10)
            cur.next = ListNode(v)
            cur = cur.next
            l1 = l1.next if l1 else None
            l2 = l2.next if l2 else None
        return dummy.next
\end{code}
\end{problem}

\begin{problem}{141}{Linked List Cycle}{Easy}
\summary{Does the list contain a cycle?}
\keyidea{Floyd's tortoise and hare: if fast ever meets slow, there is a cycle.}
\complexity{O(n) time, O(1) space}
\begin{code}
class Solution:
    def hasCycle(self, head: Optional[ListNode]) -> bool:
        slow = fast = head
        while fast and fast.next:
            slow, fast = slow.next, fast.next.next
            if slow is fast:
                return True
        return False
\end{code}
\end{problem}

\begin{problem}{287}{Find the Duplicate Number}{Medium}
\summary{Array of n+1 ints in [1, n] has one repeated value; find it without modifying the array, O(1) space.}
\keyidea{Treat i -\textgreater{} nums[i] as a linked list; the duplicate is the cycle entrance (Floyd, phase 2).}
\complexity{O(n) time, O(1) space}
\begin{code}
class Solution:
    def findDuplicate(self, nums: List[int]) -> int:
        slow = fast = 0
        while True:
            slow = nums[slow]
            fast = nums[nums[fast]]
            if slow == fast:
                break
        slow2 = 0
        while slow != slow2:
            slow, slow2 = nums[slow], nums[slow2]
        return slow
\end{code}
\end{problem}

\begin{problem}{146}{LRU Cache}{Medium}
\summary{Fixed-capacity cache with O(1) get/put that evicts the least recently used key.}
\keyidea{Hash map key -\textgreater{} node, plus a doubly linked list ordered by recency (LRU at left, MRU at right).}
\complexity{O(1) per op, O(capacity) space}
\begin{code}
class DNode:
    def __init__(self, key=0, val=0):
        self.key, self.val = key, val
        self.prev = self.next = None

class LRUCache:
    def __init__(self, capacity: int):
        self.cap, self.cache = capacity, {}
        self.left, self.right = DNode(), DNode()      # left = LRU, right = MRU
        self.left.next, self.right.prev = self.right, self.left

    def _remove(self, node):
        node.prev.next, node.next.prev = node.next, node.prev

    def _insert(self, node):                         # insert before right
        prev = self.right.prev
        prev.next = node
        node.prev, node.next = prev, self.right
        self.right.prev = node

    def get(self, key: int) -> int:
        if key not in self.cache:
            return -1
        node = self.cache[key]
        self._remove(node)
        self._insert(node)
        return node.val

    def put(self, key: int, value: int) -> None:
        if key in self.cache:
            self._remove(self.cache[key])
        self.cache[key] = DNode(key, value)
        self._insert(self.cache[key])
        if len(self.cache) > self.cap:
            lru = self.left.next
            self._remove(lru)
            del self.cache[lru.key]
\end{code}
\end{problem}

\begin{problem}{23}{Merge k Sorted Lists}{Hard}
\summary{Merge k sorted linked lists into one sorted list.}
\keyidea{Min-heap of (value, list index, node); pop the smallest and push its next.}
\complexity{O(N log k) time, O(k) space}
\begin{code}
class Solution:
    def mergeKLists(self, lists: List[Optional[ListNode]]) -> Optional[ListNode]:
        heap = [(node.val, i, node) for i, node in enumerate(lists) if node]
        heapq.heapify(heap)
        dummy = tail = ListNode()
        while heap:
            _, i, node = heapq.heappop(heap)
            tail.next = node
            tail = node
            if node.next:
                heapq.heappush(heap, (node.next.val, i, node.next))
        return dummy.next
\end{code}
\end{problem}

\begin{problem}{25}{Reverse Nodes in k-Group}{Hard}
\summary{Reverse the list k nodes at a time; a leftover tail shorter than k stays as is.}
\keyidea{For each group: find the k-th node; reverse the group; reconnect the previous group's tail.}
\complexity{O(n) time, O(1) space}
\begin{code}
class Solution:
    def reverseKGroup(self, head: Optional[ListNode], k: int) -> Optional[ListNode]:
        dummy = ListNode(0, head)
        group_prev = dummy
        while True:
            kth = group_prev
            for _ in range(k):
                kth = kth.next
                if not kth:
                    return dummy.next
            group_next = kth.next
            prev, cur = group_next, group_prev.next
            while cur is not group_next:
                nxt = cur.next
                cur.next = prev
                prev, cur = cur, nxt
            first = group_prev.next          # becomes the group's tail
            group_prev.next = kth
            group_prev = first
\end{code}
\end{problem}

\newpage
% ===== Trees (15 problems) =====
\section{Trees {\normalsize\color{muted}$\cdot$ 15 problems}}
\markboth{Trees}{}

\begin{problem}{226}{Invert Binary Tree}{Easy}
\summary{Mirror the tree (swap every left and right child).}
\keyidea{Swap children, recurse on both.}
\complexity{O(n) time, O(h) space}
\begin{code}
class Solution:
    def invertTree(self, root: Optional[TreeNode]) -> Optional[TreeNode]:
        if root:
            root.left, root.right = self.invertTree(root.right), self.invertTree(root.left)
        return root
\end{code}
\end{problem}

\begin{problem}{104}{Maximum Depth of Binary Tree}{Easy}
\summary{Number of nodes on the longest root-to-leaf path.}
\keyidea{depth = 1 + max(depth(left), depth(right)).}
\complexity{O(n) time, O(h) space}
\begin{code}
class Solution:
    def maxDepth(self, root: Optional[TreeNode]) -> int:
        if not root:
            return 0
        return 1 + max(self.maxDepth(root.left), self.maxDepth(root.right))
\end{code}
\end{problem}

\begin{problem}{543}{Diameter of Binary Tree}{Easy}
\summary{Longest path (in edges) between any two nodes.}
\keyidea{DFS returns height; at each node update best with left height + right height.}
\complexity{O(n) time, O(h) space}
\begin{code}
class Solution:
    def diameterOfBinaryTree(self, root: Optional[TreeNode]) -> int:
        best = 0
        def height(node):
            nonlocal best
            if not node:
                return 0
            l, r = height(node.left), height(node.right)
            best = max(best, l + r)
            return 1 + max(l, r)
        height(root)
        return best
\end{code}
\end{problem}

\begin{problem}{110}{Balanced Binary Tree}{Easy}
\summary{Is every node's left/right height difference at most 1?}
\keyidea{Bottom-up DFS returning height, or -1 as an 'unbalanced' sentinel.}
\complexity{O(n) time, O(h) space}
\begin{code}
class Solution:
    def isBalanced(self, root: Optional[TreeNode]) -> bool:
        def h(node):
            if not node:
                return 0
            l, r = h(node.left), h(node.right)
            if l < 0 or r < 0 or abs(l - r) > 1:
                return -1
            return 1 + max(l, r)
        return h(root) >= 0
\end{code}
\end{problem}

\begin{problem}{100}{Same Tree}{Easy}
\summary{Are two trees identical in structure and values?}
\keyidea{Both None -\textgreater{} True; one None or values differ -\textgreater{} False; else recurse on both sides.}
\complexity{O(n) time, O(h) space}
\begin{code}
class Solution:
    def isSameTree(self, p: Optional[TreeNode], q: Optional[TreeNode]) -> bool:
        if not p and not q:
            return True
        if not p or not q or p.val != q.val:
            return False
        return self.isSameTree(p.left, q.left) and self.isSameTree(p.right, q.right)
\end{code}
\end{problem}

\begin{problem}{572}{Subtree of Another Tree}{Easy}
\summary{Is subRoot an exact subtree somewhere inside root?}
\keyidea{At every node of root, test sameTree(node, subRoot).}
\complexity{O(m * n) time, O(h) space}
\begin{code}
class Solution:
    def isSubtree(self, root: Optional[TreeNode], subRoot: Optional[TreeNode]) -> bool:
        def same(a, b):
            if not a and not b:
                return True
            if not a or not b or a.val != b.val:
                return False
            return same(a.left, b.left) and same(a.right, b.right)
        if not subRoot:
            return True
        if not root:
            return False
        return (same(root, subRoot) or self.isSubtree(root.left, subRoot)
                or self.isSubtree(root.right, subRoot))
\end{code}
\end{problem}

\begin{problem}{235}{Lowest Common Ancestor of a Binary Search Tree}{Medium}
\summary{Find the LCA of nodes p and q in a BST.}
\keyidea{Both smaller -\textgreater{} go left; both larger -\textgreater{} go right; otherwise this node is where they split.}
\complexity{O(h) time, O(1) space}
\begin{code}
class Solution:
    def lowestCommonAncestor(self, root: TreeNode, p: TreeNode, q: TreeNode) -> TreeNode:
        cur = root
        while cur:
            if p.val < cur.val and q.val < cur.val:
                cur = cur.left
            elif p.val > cur.val and q.val > cur.val:
                cur = cur.right
            else:
                return cur
\end{code}
\end{problem}

\begin{problem}{102}{Binary Tree Level Order Traversal}{Medium}
\summary{Return node values level by level.}
\keyidea{BFS with a queue; process len(queue) nodes per level.}
\complexity{O(n) time, O(n) space}
\begin{code}
class Solution:
    def levelOrder(self, root: Optional[TreeNode]) -> List[List[int]]:
        res, q = [], deque([root] if root else [])
        while q:
            level = []
            for _ in range(len(q)):
                node = q.popleft()
                level.append(node.val)
                if node.left: q.append(node.left)
                if node.right: q.append(node.right)
            res.append(level)
        return res
\end{code}
\end{problem}

\begin{problem}{199}{Binary Tree Right Side View}{Medium}
\summary{Values visible when looking at the tree from the right, top to bottom.}
\keyidea{BFS by level; take the last node of each level.}
\complexity{O(n) time, O(n) space}
\begin{code}
class Solution:
    def rightSideView(self, root: Optional[TreeNode]) -> List[int]:
        res, q = [], deque([root] if root else [])
        while q:
            for i in range(len(q)):
                node = q.popleft()
                if i == 0:
                    res.append(node.val)       # enqueue right first -> first is rightmost
                if node.right: q.append(node.right)
                if node.left: q.append(node.left)
        return res
\end{code}
\end{problem}

\begin{problem}{1448}{Count Good Nodes in Binary Tree}{Medium}
\summary{A node is good if no node on the path from the root to it is larger. Count good nodes.}
\keyidea{DFS carrying the max value seen on the path.}
\complexity{O(n) time, O(h) space}
\begin{code}
class Solution:
    def goodNodes(self, root: TreeNode) -> int:
        def dfs(node, mx):
            if not node:
                return 0
            good = 1 if node.val >= mx else 0
            mx = max(mx, node.val)
            return good + dfs(node.left, mx) + dfs(node.right, mx)
        return dfs(root, root.val)
\end{code}
\end{problem}

\begin{problem}{98}{Validate Binary Search Tree}{Medium}
\summary{Is the tree a valid BST (strict ordering)?}
\keyidea{DFS passing an open interval (low, high) that every node must sit inside.}
\complexity{O(n) time, O(h) space}
\begin{code}
class Solution:
    def isValidBST(self, root: Optional[TreeNode]) -> bool:
        def ok(node, lo, hi):
            if not node:
                return True
            if not (lo < node.val < hi):
                return False
            return ok(node.left, lo, node.val) and ok(node.right, node.val, hi)
        return ok(root, float("-inf"), float("inf"))
\end{code}
\end{problem}

\begin{problem}{230}{Kth Smallest Element in a BST}{Medium}
\summary{Return the k-th smallest value (1-indexed) in a BST.}
\keyidea{In-order traversal visits values in sorted order; stop at the k-th pop.}
\complexity{O(h + k) time, O(h) space}
\begin{code}
class Solution:
    def kthSmallest(self, root: Optional[TreeNode], k: int) -> int:
        st, cur = [], root
        while st or cur:
            while cur:
                st.append(cur)
                cur = cur.left
            cur = st.pop()
            k -= 1
            if k == 0:
                return cur.val
            cur = cur.right
\end{code}
\end{problem}

\begin{problem}{105}{Construct Binary Tree from Preorder and Inorder Traversal}{Medium}
\summary{Rebuild the tree from its preorder and inorder sequences (unique values).}
\keyidea{preorder[0] is the root; its index in inorder splits left/right subtrees. Use a value-\textgreater{}index map.}
\complexity{O(n) time, O(n) space}
\begin{code}
class Solution:
    def buildTree(self, preorder: List[int], inorder: List[int]) -> Optional[TreeNode]:
        idx = {v: i for i, v in enumerate(inorder)}
        pre = iter(preorder)
        def build(l, r):
            if l > r:
                return None
            val = next(pre)
            node = TreeNode(val)
            m = idx[val]
            node.left = build(l, m - 1)
            node.right = build(m + 1, r)
            return node
        return build(0, len(inorder) - 1)
\end{code}
\end{problem}

\begin{problem}{124}{Binary Tree Maximum Path Sum}{Hard}
\summary{Max sum of any path (any start and end node, must be connected).}
\keyidea{DFS returns the best downward 'arm' (clamped at 0); at each node update answer with val + left + right.}
\complexity{O(n) time, O(h) space}
\begin{code}
class Solution:
    def maxPathSum(self, root: Optional[TreeNode]) -> int:
        best = root.val
        def arm(node):
            nonlocal best
            if not node:
                return 0
            l = max(arm(node.left), 0)
            r = max(arm(node.right), 0)
            best = max(best, node.val + l + r)
            return node.val + max(l, r)
        arm(root)
        return best
\end{code}
\end{problem}

\begin{problem}{297}{Serialize and Deserialize Binary Tree}{Hard}
\summary{Convert a tree to a string and back.}
\keyidea{Preorder with 'N' for nulls; deserialize by consuming tokens in the same preorder.}
\complexity{O(n) time, O(n) space}
\begin{code}
class Codec:
    def serialize(self, root):
        out = []
        def dfs(node):
            if not node:
                out.append("N")
                return
            out.append(str(node.val))
            dfs(node.left)
            dfs(node.right)
        dfs(root)
        return ",".join(out)

    def deserialize(self, data):
        vals = iter(data.split(","))
        def dfs():
            v = next(vals)
            if v == "N":
                return None
            node = TreeNode(int(v))
            node.left = dfs()
            node.right = dfs()
            return node
        return dfs()
\end{code}
\end{problem}

\newpage
% ===== Tries (3 problems) =====
\section{Tries \texorpdfstring{{\normalsize\color{muted}$\cdot$ 3 problems}}{- 3 problems}}
\markboth{Tries}{}

\begin{problem}{208}{Implement Trie (Prefix Tree)}{Medium}
\summary{Build insert, search (whole word) and startsWith (prefix).}
\keyidea{Each node: dict of children + end-of-word flag.}
\complexity{O(L) per op, O(total chars) space}
\begin{code}
class TrieNode:
    def __init__(self):
        self.children = {}
        self.end = False

class Trie:
    def __init__(self):
        self.root = TrieNode()

    def insert(self, word: str) -> None:
        cur = self.root
        for c in word:
            cur = cur.children.setdefault(c, TrieNode())
        cur.end = True

    def _walk(self, s):
        cur = self.root
        for c in s:
            if c not in cur.children:
                return None
            cur = cur.children[c]
        return cur

    def search(self, word: str) -> bool:
        node = self._walk(word)
        return bool(node and node.end)

    def startsWith(self, prefix: str) -> bool:
        return self._walk(prefix) is not None
\end{code}
\end{problem}

\begin{problem}{211}{Design Add and Search Words Data Structure}{Medium}
\summary{Like a trie, but search may contain '.' matching any one letter.}
\keyidea{Trie + DFS; on '.', try every child.}
\complexity{add O(L); search O(26\textasciicircum{}dots * L) worst case}
\begin{code}
class TrieNode:
    def __init__(self):
        self.children = {}
        self.end = False

class WordDictionary:
    def __init__(self):
        self.root = TrieNode()

    def addWord(self, word: str) -> None:
        cur = self.root
        for c in word:
            cur = cur.children.setdefault(c, TrieNode())
        cur.end = True

    def search(self, word: str) -> bool:
        def dfs(i, node):
            if i == len(word):
                return node.end
            c = word[i]
            if c == ".":
                return any(dfs(i + 1, ch) for ch in node.children.values())
            return c in node.children and dfs(i + 1, node.children[c])
        return dfs(0, self.root)
\end{code}
\end{problem}

\begin{problem}{212}{Word Search II}{Hard}
\summary{Find all dictionary words that can be traced on a letter grid (adjacent cells, no reuse).}
\keyidea{Insert words into a trie; DFS from each cell following trie edges; prune found words.}
\complexity{O(m * n * 4\textasciicircum{}L) worst case, O(total word chars) space}
\begin{code}
class Solution:
    def findWords(self, board: List[List[str]], words: List[str]) -> List[str]:
        root = {}
        for w in words:
            node = root
            for c in w:
                node = node.setdefault(c, {})
            node["$"] = w
        R, C = len(board), len(board[0])
        res = []
        def dfs(r, c, parent):
            ch = board[r][c]
            node = parent[ch]
            if "$" in node:
                res.append(node.pop("$"))     # avoid duplicates
            board[r][c] = "#"
            for dr, dc in ((1, 0), (-1, 0), (0, 1), (0, -1)):
                nr, nc = r + dr, c + dc
                if 0 <= nr < R and 0 <= nc < C and board[nr][nc] in node:
                    dfs(nr, nc, node)
            board[r][c] = ch
            if not node:                       # prune empty branch
                parent.pop(ch)
        for r in range(R):
            for c in range(C):
                if board[r][c] in root:
                    dfs(r, c, root)
        return res
\end{code}
\end{problem}

\newpage
% ===== Heap / Priority Queue (7 problems) =====
\section{Heap / Priority Queue \texorpdfstring{{\normalsize\color{muted}$\cdot$ 7 problems}}{- 7 problems}}
\markboth{Heap / Priority Queue}{}

\begin{problem}{703}{Kth Largest Element in a Stream}{Easy}
\summary{Class that returns the k-th largest after each added value.}
\keyidea{Keep a min-heap of size k; its top is the k-th largest.}
\complexity{O(log k) per add, O(k) space}
\begin{code}
class KthLargest:
    def __init__(self, k: int, nums: List[int]):
        self.k, self.heap = k, nums
        heapq.heapify(self.heap)
        while len(self.heap) > k:
            heapq.heappop(self.heap)

    def add(self, val: int) -> int:
        heapq.heappush(self.heap, val)
        if len(self.heap) > self.k:
            heapq.heappop(self.heap)
        return self.heap[0]
\end{code}
\end{problem}

\begin{problem}{1046}{Last Stone Weight}{Easy}
\summary{Repeatedly smash the two heaviest stones; return the last stone's weight (or 0).}
\keyidea{Max-heap via negated values.}
\complexity{O(n log n) time, O(n) space}
\begin{code}
class Solution:
    def lastStoneWeight(self, stones: List[int]) -> int:
        h = [-s for s in stones]
        heapq.heapify(h)
        while len(h) > 1:
            a, b = -heapq.heappop(h), -heapq.heappop(h)
            if a != b:
                heapq.heappush(h, -(a - b))
        return -h[0] if h else 0
\end{code}
\end{problem}

\begin{problem}{973}{K Closest Points to Origin}{Medium}
\summary{Return the k points nearest to (0, 0).}
\keyidea{Max-heap of size k on squared distance (negated); drop the farthest when over k.}
\complexity{O(n log k) time, O(k) space}
\begin{code}
class Solution:
    def kClosest(self, points: List[List[int]], k: int) -> List[List[int]]:
        h = []
        for x, y in points:
            heapq.heappush(h, (-(x * x + y * y), x, y))
            if len(h) > k:
                heapq.heappop(h)
        return [[x, y] for _, x, y in h]
\end{code}
\end{problem}

\begin{problem}{215}{Kth Largest Element in an Array}{Medium}
\summary{Return the k-th largest value in an unsorted array.}
\keyidea{Min-heap of size k (or Quickselect for average O(n)).}
\complexity{O(n log k) time, O(k) space}
\begin{code}
class Solution:
    def findKthLargest(self, nums: List[int], k: int) -> int:
        h = []
        for n in nums:
            heapq.heappush(h, n)
            if len(h) > k:
                heapq.heappop(h)
        return h[0]
\end{code}
\end{problem}

\begin{problem}{621}{Task Scheduler}{Medium}
\summary{Same tasks need n idle slots between them; minimum total time units to run all.}
\keyidea{Formula: (maxCount - 1) * (n + 1) + numberOfTasksWithMaxCount, but at least len(tasks).}
\complexity{O(N) time, O(26) space}
\begin{code}
class Solution:
    def leastInterval(self, tasks: List[str], n: int) -> int:
        count = Counter(tasks)
        mx = max(count.values())
        num_mx = sum(1 for v in count.values() if v == mx)
        return max(len(tasks), (mx - 1) * (n + 1) + num_mx)
\end{code}
\end{problem}

\begin{problem}{355}{Design Twitter}{Medium}
\summary{postTweet, follow, unfollow, getNewsFeed (10 most recent from self + followees).}
\keyidea{Global decreasing time counter; merge the followees' tweet lists with a heap (k-way merge).}
\complexity{getNewsFeed O(F log F) for F followees; O(users + tweets) space}
\begin{code}
class Twitter:
    def __init__(self):
        self.time = 0
        self.tweets = defaultdict(list)    # user -> [(time, tweetId)]
        self.following = defaultdict(set)

    def postTweet(self, userId: int, tweetId: int) -> None:
        self.tweets[userId].append((self.time, tweetId))
        self.time -= 1                     # smaller = newer (min-heap)

    def getNewsFeed(self, userId: int) -> List[int]:
        res, h = [], []
        users = self.following[userId] | {userId}
        for u in users:
            if self.tweets[u]:
                i = len(self.tweets[u]) - 1
                t, tid = self.tweets[u][i]
                h.append((t, tid, u, i - 1))
        heapq.heapify(h)
        while h and len(res) < 10:
            t, tid, u, i = heapq.heappop(h)
            res.append(tid)
            if i >= 0:
                t2, tid2 = self.tweets[u][i]
                heapq.heappush(h, (t2, tid2, u, i - 1))
        return res

    def follow(self, followerId: int, followeeId: int) -> None:
        if followerId != followeeId:
            self.following[followerId].add(followeeId)

    def unfollow(self, followerId: int, followeeId: int) -> None:
        self.following[followerId].discard(followeeId)
\end{code}
\end{problem}

\begin{problem}{295}{Find Median from Data Stream}{Hard}
\summary{Support addNum and findMedian on a growing stream.}
\keyidea{Max-heap for the lower half, min-heap for the upper half; keep sizes within 1.}
\complexity{addNum O(log n), findMedian O(1)}
\begin{code}
class MedianFinder:
    def __init__(self):
        self.small, self.large = [], []    # small = max-heap (negated)

    def addNum(self, num: int) -> None:
        heapq.heappush(self.small, -num)
        heapq.heappush(self.large, -heapq.heappop(self.small))
        if len(self.large) > len(self.small):
            heapq.heappush(self.small, -heapq.heappop(self.large))

    def findMedian(self) -> float:
        if len(self.small) > len(self.large):
            return -self.small[0]
        return (-self.small[0] + self.large[0]) / 2
\end{code}
\end{problem}

\newpage
% ===== Backtracking (9 problems) =====
\section{Backtracking \texorpdfstring{{\normalsize\color{muted}$\cdot$ 9 problems}}{- 9 problems}}
\markboth{Backtracking}{}

\begin{problem}{78}{Subsets}{Medium}
\summary{Return all subsets of distinct integers.}
\keyidea{At each index choose: include nums[i] or skip it.}
\complexity{O(n * 2\textasciicircum{}n) time, O(n) recursion}
\begin{code}
class Solution:
    def subsets(self, nums: List[int]) -> List[List[int]]:
        res, cur = [], []
        def dfs(i):
            if i == len(nums):
                res.append(cur[:])
                return
            cur.append(nums[i]); dfs(i + 1); cur.pop()
            dfs(i + 1)
        dfs(0)
        return res
\end{code}
\end{problem}

\begin{problem}{39}{Combination Sum}{Medium}
\summary{Unique combinations of candidates summing to target; each candidate may be reused.}
\keyidea{Branch: use candidates[i] again (stay at i) or move to i+1. Stop when total \textgreater{}= target.}
\complexity{O(2\textasciicircum{}(t/min)) time, O(t/min) recursion}
\begin{code}
class Solution:
    def combinationSum(self, candidates: List[int], target: int) -> List[List[int]]:
        res, cur = [], []
        def dfs(i, total):
            if total == target:
                res.append(cur[:])
                return
            if i == len(candidates) or total > target:
                return
            cur.append(candidates[i]); dfs(i, total + candidates[i]); cur.pop()
            dfs(i + 1, total)
        dfs(0, 0)
        return res
\end{code}
\end{problem}

\begin{problem}{46}{Permutations}{Medium}
\summary{Return all orderings of distinct integers.}
\keyidea{Build position by position, using each unused number once.}
\complexity{O(n * n!) time, O(n) recursion}
\begin{code}
class Solution:
    def permute(self, nums: List[int]) -> List[List[int]]:
        res, cur, used = [], [], [False] * len(nums)
        def dfs():
            if len(cur) == len(nums):
                res.append(cur[:])
                return
            for i, n in enumerate(nums):
                if not used[i]:
                    used[i] = True
                    cur.append(n); dfs(); cur.pop()
                    used[i] = False
        dfs()
        return res
\end{code}
\end{problem}

\begin{problem}{90}{Subsets II}{Medium}
\summary{All unique subsets when the input may contain duplicates.}
\keyidea{Sort; after choosing to skip nums[i], also skip all equal values.}
\complexity{O(n * 2\textasciicircum{}n) time, O(n) recursion}
\begin{code}
class Solution:
    def subsetsWithDup(self, nums: List[int]) -> List[List[int]]:
        nums.sort()
        res, cur = [], []
        def dfs(i):
            if i == len(nums):
                res.append(cur[:])
                return
            cur.append(nums[i]); dfs(i + 1); cur.pop()
            while i + 1 < len(nums) and nums[i] == nums[i + 1]:
                i += 1
            dfs(i + 1)
        dfs(0)
        return res
\end{code}
\end{problem}

\begin{problem}{40}{Combination Sum II}{Medium}
\summary{Each candidate used at most once, input has duplicates; unique combos summing to target.}
\keyidea{Sort; loop over choices from i, skipping j \textgreater{} i where c[j] == c[j-1]; break when too big.}
\complexity{O(2\textasciicircum{}n) time, O(n) recursion}
\begin{code}
class Solution:
    def combinationSum2(self, candidates: List[int], target: int) -> List[List[int]]:
        candidates.sort()
        res, cur = [], []
        def dfs(i, remain):
            if remain == 0:
                res.append(cur[:])
                return
            for j in range(i, len(candidates)):
                if j > i and candidates[j] == candidates[j - 1]:
                    continue
                if candidates[j] > remain:
                    break
                cur.append(candidates[j]); dfs(j + 1, remain - candidates[j]); cur.pop()
        dfs(0, target)
        return res
\end{code}
\end{problem}

\begin{problem}{79}{Word Search}{Medium}
\summary{Can word be traced through adjacent grid cells, each cell used at most once?}
\keyidea{DFS from each cell; mark visited in place, restore on backtrack.}
\complexity{O(m * n * 4\textasciicircum{}L) time, O(L) recursion}
\begin{code}
class Solution:
    def exist(self, board: List[List[str]], word: str) -> bool:
        R, C = len(board), len(board[0])
        def dfs(r, c, i):
            if i == len(word):
                return True
            if not (0 <= r < R and 0 <= c < C) or board[r][c] != word[i]:
                return False
            tmp, board[r][c] = board[r][c], "#"
            found = (dfs(r + 1, c, i + 1) or dfs(r - 1, c, i + 1) or
                     dfs(r, c + 1, i + 1) or dfs(r, c - 1, i + 1))
            board[r][c] = tmp
            return found
        return any(dfs(r, c, 0) for r in range(R) for c in range(C))
\end{code}
\end{problem}

\begin{problem}{131}{Palindrome Partitioning}{Medium}
\summary{Split s every possible way so that every piece is a palindrome.}
\keyidea{From index i, try every end j where s[i:j+1] is a palindrome; recurse from j+1.}
\complexity{O(n * 2\textasciicircum{}n) time, O(n) recursion}
\begin{code}
class Solution:
    def partition(self, s: str) -> List[List[str]]:
        res, cur = [], []
        def dfs(i):
            if i == len(s):
                res.append(cur[:])
                return
            for j in range(i, len(s)):
                piece = s[i : j + 1]
                if piece == piece[::-1]:
                    cur.append(piece); dfs(j + 1); cur.pop()
        dfs(0)
        return res
\end{code}
\end{problem}

\begin{problem}{17}{Letter Combinations of a Phone Number}{Medium}
\summary{All letter strings a digit string (2-9) could represent on a phone keypad.}
\keyidea{Backtrack over digits, appending each mapped letter.}
\complexity{O(n * 4\textasciicircum{}n) time, O(n) recursion}
\begin{code}
class Solution:
    def letterCombinations(self, digits: str) -> List[str]:
        if not digits:
            return []
        keys = {"2": "abc", "3": "def", "4": "ghi", "5": "jkl",
                "6": "mno", "7": "pqrs", "8": "tuv", "9": "wxyz"}
        res = []
        def dfs(i, cur):
            if i == len(digits):
                res.append(cur)
                return
            for ch in keys[digits[i]]:
                dfs(i + 1, cur + ch)
        dfs(0, "")
        return res
\end{code}
\end{problem}

\begin{problem}{51}{N-Queens}{Hard}
\summary{All ways to place n queens on an n x n board with no two attacking.}
\keyidea{Row by row; track used columns, (r + c) diagonals and (r - c) anti-diagonals in sets.}
\complexity{O(n!) time, O(n\textasciicircum{}2) for board}
\begin{code}
class Solution:
    def solveNQueens(self, n: int) -> List[List[str]]:
        cols, pos, neg = set(), set(), set()
        board = [["."] * n for _ in range(n)]
        res = []
        def dfs(r):
            if r == n:
                res.append(["".join(row) for row in board])
                return
            for c in range(n):
                if c in cols or (r + c) in pos or (r - c) in neg:
                    continue
                cols.add(c); pos.add(r + c); neg.add(r - c); board[r][c] = "Q"
                dfs(r + 1)
                cols.remove(c); pos.remove(r + c); neg.remove(r - c); board[r][c] = "."
        dfs(0)
        return res
\end{code}
\end{problem}

\newpage
% ===== Graphs (13 problems) =====
\section{Graphs \texorpdfstring{{\normalsize\color{muted}$\cdot$ 13 problems}}{- 13 problems}}
\markboth{Graphs}{}

\begin{problem}{200}{Number of Islands}{Medium}
\summary{Count groups of connected '1's (4-directional) in a grid.}
\keyidea{Scan; at each unvisited land cell, count +1 and flood-fill (DFS/BFS) the whole island.}
\complexity{O(m * n) time and space}
\begin{code}
class Solution:
    def numIslands(self, grid: List[List[str]]) -> int:
        R, C = len(grid), len(grid[0])
        def sink(r, c):
            if 0 <= r < R and 0 <= c < C and grid[r][c] == "1":
                grid[r][c] = "0"
                sink(r + 1, c); sink(r - 1, c); sink(r, c + 1); sink(r, c - 1)
        count = 0
        for r in range(R):
            for c in range(C):
                if grid[r][c] == "1":
                    count += 1
                    sink(r, c)
        return count
\end{code}
\end{problem}

\begin{problem}{133}{Clone Graph}{Medium}
\summary{Deep-copy a connected undirected graph given one node.}
\keyidea{DFS with a map old -\textgreater{} clone; create the clone before visiting neighbours (handles cycles).}
\complexity{O(V + E) time, O(V) space}
\begin{code}
class Solution:
    def cloneGraph(self, node: Optional["Node"]) -> Optional["Node"]:
        clones = {}
        def dfs(n):
            if n in clones:
                return clones[n]
            copy = Node(n.val)
            clones[n] = copy
            for nb in n.neighbors:
                copy.neighbors.append(dfs(nb))
            return copy
        return dfs(node) if node else None
\end{code}
\end{problem}

\begin{problem}{695}{Max Area of Island}{Medium}
\summary{Size of the largest island of 1s.}
\keyidea{Flood-fill returning the area; take the max.}
\complexity{O(m * n) time and space}
\begin{code}
class Solution:
    def maxAreaOfIsland(self, grid: List[List[int]]) -> int:
        R, C = len(grid), len(grid[0])
        def area(r, c):
            if not (0 <= r < R and 0 <= c < C) or grid[r][c] != 1:
                return 0
            grid[r][c] = 0
            return 1 + area(r + 1, c) + area(r - 1, c) + area(r, c + 1) + area(r, c - 1)
        return max(area(r, c) for r in range(R) for c in range(C))
\end{code}
\end{problem}

\begin{problem}{417}{Pacific Atlantic Water Flow}{Medium}
\summary{Cells whose water can flow (to equal or lower neighbours) to both the Pacific (top/left) and Atlantic (bottom/right).}
\keyidea{Reverse it: DFS uphill from each ocean's border; answer = intersection of both reachable sets.}
\complexity{O(m * n) time and space}
\begin{code}
class Solution:
    def pacificAtlantic(self, heights: List[List[int]]) -> List[List[int]]:
        R, C = len(heights), len(heights[0])
        pac, atl = set(), set()
        def dfs(r, c, seen, prev):
            if ((r, c) in seen or not (0 <= r < R and 0 <= c < C)
                    or heights[r][c] < prev):
                return
            seen.add((r, c))
            for dr, dc in ((1, 0), (-1, 0), (0, 1), (0, -1)):
                dfs(r + dr, c + dc, seen, heights[r][c])
        for c in range(C):
            dfs(0, c, pac, heights[0][c])
            dfs(R - 1, c, atl, heights[R - 1][c])
        for r in range(R):
            dfs(r, 0, pac, heights[r][0])
            dfs(r, C - 1, atl, heights[r][C - 1])
        return [list(p) for p in pac & atl]
\end{code}
\end{problem}

\begin{problem}{130}{Surrounded Regions}{Medium}
\summary{Flip every 'O' region not connected to the border into 'X'.}
\keyidea{Mark border-connected O's as safe (T) by DFS; flip the rest to X; restore T to O.}
\complexity{O(m * n) time and space}
\begin{code}
class Solution:
    def solve(self, board: List[List[str]]) -> None:
        R, C = len(board), len(board[0])
        def mark(r, c):
            if 0 <= r < R and 0 <= c < C and board[r][c] == "O":
                board[r][c] = "T"
                mark(r + 1, c); mark(r - 1, c); mark(r, c + 1); mark(r, c - 1)
        for r in range(R):
            mark(r, 0); mark(r, C - 1)
        for c in range(C):
            mark(0, c); mark(R - 1, c)
        for r in range(R):
            for c in range(C):
                if board[r][c] == "O":
                    board[r][c] = "X"
                elif board[r][c] == "T":
                    board[r][c] = "O"
\end{code}
\end{problem}

\begin{problem}{994}{Rotting Oranges}{Medium}
\summary{Each minute rot spreads to adjacent fresh oranges. Minutes until none are fresh (-1 if impossible).}
\keyidea{Multi-source BFS from all rotten oranges at once; count levels.}
\complexity{O(m * n) time and space}
\begin{code}
class Solution:
    def orangesRotting(self, grid: List[List[int]]) -> int:
        R, C = len(grid), len(grid[0])
        q, fresh = deque(), 0
        for r in range(R):
            for c in range(C):
                if grid[r][c] == 2: q.append((r, c))
                elif grid[r][c] == 1: fresh += 1
        minutes = 0
        while q and fresh:
            for _ in range(len(q)):
                r, c = q.popleft()
                for nr, nc in ((r + 1, c), (r - 1, c), (r, c + 1), (r, c - 1)):
                    if 0 <= nr < R and 0 <= nc < C and grid[nr][nc] == 1:
                        grid[nr][nc] = 2
                        fresh -= 1
                        q.append((nr, nc))
            minutes += 1
        return minutes if fresh == 0 else -1
\end{code}
\end{problem}

\begin{problem}{286}{Walls and Gates}{Medium}
\summary{Fill each empty room (INF) with its distance to the nearest gate (0); walls are -1.}
\keyidea{Multi-source BFS starting from every gate.}
\complexity{O(m * n) time and space}
\begin{code}
class Solution:
    def wallsAndGates(self, rooms: List[List[int]]) -> None:
        INF = 2147483647
        R, C = len(rooms), len(rooms[0])
        q = deque((r, c) for r in range(R) for c in range(C) if rooms[r][c] == 0)
        while q:
            r, c = q.popleft()
            for nr, nc in ((r + 1, c), (r - 1, c), (r, c + 1), (r, c - 1)):
                if 0 <= nr < R and 0 <= nc < C and rooms[nr][nc] == INF:
                    rooms[nr][nc] = rooms[r][c] + 1
                    q.append((nr, nc))
\end{code}
\end{problem}

\begin{problem}{207}{Course Schedule}{Medium}
\summary{Given prerequisite pairs, can all courses be finished (i.e. no cycle)?}
\keyidea{Kahn's topological sort: repeatedly take in-degree-0 nodes; success if all are taken.}
\complexity{O(V + E) time and space}
\begin{code}
class Solution:
    def canFinish(self, numCourses: int, prerequisites: List[List[int]]) -> bool:
        adj = [[] for _ in range(numCourses)]
        indeg = [0] * numCourses
        for a, b in prerequisites:          # b -> a
            adj[b].append(a)
            indeg[a] += 1
        q = deque(i for i in range(numCourses) if indeg[i] == 0)
        taken = 0
        while q:
            c = q.popleft()
            taken += 1
            for nxt in adj[c]:
                indeg[nxt] -= 1
                if indeg[nxt] == 0:
                    q.append(nxt)
        return taken == numCourses
\end{code}
\end{problem}

\begin{problem}{210}{Course Schedule II}{Medium}
\summary{Return a valid order to take all courses, or [] if impossible.}
\keyidea{Same Kahn's algorithm; the pop order is the answer.}
\complexity{O(V + E) time and space}
\begin{code}
class Solution:
    def findOrder(self, numCourses: int, prerequisites: List[List[int]]) -> List[int]:
        adj = [[] for _ in range(numCourses)]
        indeg = [0] * numCourses
        for a, b in prerequisites:
            adj[b].append(a)
            indeg[a] += 1
        q = deque(i for i in range(numCourses) if indeg[i] == 0)
        order = []
        while q:
            c = q.popleft()
            order.append(c)
            for nxt in adj[c]:
                indeg[nxt] -= 1
                if indeg[nxt] == 0:
                    q.append(nxt)
        return order if len(order) == numCourses else []
\end{code}
\end{problem}

\begin{problem}{684}{Redundant Connection}{Medium}
\summary{A tree plus one extra edge; return the edge (latest in input) that creates the cycle.}
\keyidea{Union-Find: the first edge whose endpoints already share a root is redundant.}
\complexity{O(n * alpha(n)) time, O(n) space}
\begin{code}
class Solution:
    def findRedundantConnection(self, edges: List[List[int]]) -> List[int]:
        parent = list(range(len(edges) + 1))
        def find(x):
            while parent[x] != x:
                parent[x] = parent[parent[x]]   # path compression
                x = parent[x]
            return x
        for a, b in edges:
            ra, rb = find(a), find(b)
            if ra == rb:
                return [a, b]
            parent[ra] = rb
\end{code}
\end{problem}

\begin{problem}{323}{Number of Connected Components in an Undirected Graph}{Medium}
\summary{n nodes and an edge list; count connected components.}
\keyidea{Union-Find: start with n components, subtract 1 for every successful union.}
\complexity{O(E * alpha(n)) time, O(n) space}
\begin{code}
class Solution:
    def countComponents(self, n: int, edges: List[List[int]]) -> int:
        parent = list(range(n))
        def find(x):
            while parent[x] != x:
                parent[x] = parent[parent[x]]
                x = parent[x]
            return x
        comps = n
        for a, b in edges:
            ra, rb = find(a), find(b)
            if ra != rb:
                parent[ra] = rb
                comps -= 1
        return comps
\end{code}
\end{problem}

\begin{problem}{261}{Graph Valid Tree}{Medium}
\summary{Do n nodes and the given undirected edges form a valid tree?}
\keyidea{A tree has exactly n-1 edges and no cycle (equivalently: n-1 edges and connected).}
\complexity{O(n * alpha(n)) time, O(n) space}
\begin{code}
class Solution:
    def validTree(self, n: int, edges: List[List[int]]) -> bool:
        if len(edges) != n - 1:
            return False
        parent = list(range(n))
        def find(x):
            while parent[x] != x:
                parent[x] = parent[parent[x]]
                x = parent[x]
            return x
        for a, b in edges:
            ra, rb = find(a), find(b)
            if ra == rb:
                return False
            parent[ra] = rb
        return True
\end{code}
\end{problem}

\begin{problem}{127}{Word Ladder}{Hard}
\summary{Fewest words in a transformation from beginWord to endWord, changing one letter at a time via wordList.}
\keyidea{BFS; neighbours found via wildcard patterns like 'h*t' -\textgreater{} words.}
\complexity{O(N * L\textasciicircum{}2) time, O(N * L\textasciicircum{}2) space}
\begin{code}
class Solution:
    def ladderLength(self, beginWord: str, endWord: str, wordList: List[str]) -> int:
        if endWord not in wordList:
            return 0
        patterns = defaultdict(list)
        for w in wordList + [beginWord]:
            for i in range(len(w)):
                patterns[w[:i] + "*" + w[i + 1:]].append(w)
        seen, q, steps = {beginWord}, deque([beginWord]), 1
        while q:
            for _ in range(len(q)):
                w = q.popleft()
                if w == endWord:
                    return steps
                for i in range(len(w)):
                    for nb in patterns[w[:i] + "*" + w[i + 1:]]:
                        if nb not in seen:
                            seen.add(nb)
                            q.append(nb)
            steps += 1
        return 0
\end{code}
\end{problem}

\newpage
% ===== Advanced Graphs (6 problems) =====
\section{Advanced Graphs {\normalsize\color{muted}$\cdot$ 6 problems}}
\markboth{Advanced Graphs}{}

\begin{problem}{332}{Reconstruct Itinerary}{Hard}
\summary{Use every ticket exactly once starting at JFK; return the lexically smallest valid route.}
\keyidea{Hierholzer's algorithm: adjacency lists sorted in reverse (pop smallest), DFS, append on backtrack, reverse at end.}
\complexity{O(E log E) time, O(E) space}
\begin{code}
class Solution:
    def findItinerary(self, tickets: List[List[str]]) -> List[str]:
        adj = defaultdict(list)
        for src, dst in sorted(tickets, reverse=True):
            adj[src].append(dst)            # pop() gives smallest
        route, st = [], ["JFK"]
        while st:
            while adj[st[-1]]:
                st.append(adj[st[-1]].pop())
            route.append(st.pop())
        return route[::-1]
\end{code}
\end{problem}

\begin{problem}{1584}{Min Cost to Connect All Points}{Medium}
\summary{Connect all points with minimum total Manhattan-distance edges.}
\keyidea{Prim's MST with a min-heap starting from point 0.}
\complexity{O(n\textasciicircum{}2 log n) time, O(n\textasciicircum{}2) space}
\begin{code}
class Solution:
    def minCostConnectPoints(self, points: List[List[int]]) -> int:
        n = len(points)
        seen, h, total = set(), [(0, 0)], 0   # (cost, point)
        while len(seen) < n:
            cost, i = heapq.heappop(h)
            if i in seen:
                continue
            seen.add(i)
            total += cost
            xi, yi = points[i]
            for j in range(n):
                if j not in seen:
                    xj, yj = points[j]
                    heapq.heappush(h, (abs(xi - xj) + abs(yi - yj), j))
        return total
\end{code}
\end{problem}

\begin{problem}{743}{Network Delay Time}{Medium}
\summary{Signal sent from node k along weighted directed edges; time for all nodes to receive it (-1 if not all).}
\keyidea{Dijkstra from k; answer is the largest shortest-distance.}
\complexity{O(E log V) time, O(V + E) space}
\begin{code}
class Solution:
    def networkDelayTime(self, times: List[List[int]], n: int, k: int) -> int:
        adj = defaultdict(list)
        for u, v, w in times:
            adj[u].append((v, w))
        dist, h = {}, [(0, k)]
        while h:
            d, u = heapq.heappop(h)
            if u in dist:
                continue
            dist[u] = d
            for v, w in adj[u]:
                if v not in dist:
                    heapq.heappush(h, (d + w, v))
        return max(dist.values()) if len(dist) == n else -1
\end{code}
\end{problem}

\begin{problem}{778}{Swim in Rising Water}{Hard}
\summary{At time t you can swim through cells with height \textless{}= t. Earliest time to go from top-left to bottom-right.}
\keyidea{Dijkstra-like: min-heap keyed by the max height seen on the path so far.}
\complexity{O(n\textasciicircum{}2 log n) time, O(n\textasciicircum{}2) space}
\begin{code}
class Solution:
    def swimInWater(self, grid: List[List[int]]) -> int:
        n = len(grid)
        h, seen = [(grid[0][0], 0, 0)], {(0, 0)}
        while h:
            t, r, c = heapq.heappop(h)
            if r == c == n - 1:
                return t
            for nr, nc in ((r + 1, c), (r - 1, c), (r, c + 1), (r, c - 1)):
                if 0 <= nr < n and 0 <= nc < n and (nr, nc) not in seen:
                    seen.add((nr, nc))
                    heapq.heappush(h, (max(t, grid[nr][nc]), nr, nc))
\end{code}
\end{problem}

\begin{problem}{269}{Alien Dictionary}{Hard}
\summary{Words are sorted in an unknown alphabet; return a valid letter order, or '' if inconsistent.}
\keyidea{Compare adjacent words: first differing letter gives an edge. Prefix trap: longer word before its prefix -\textgreater{} ''. Then topological sort.}
\complexity{O(total chars) time, O(alphabet\textasciicircum{}2) space}
\begin{code}
class Solution:
    def alienOrder(self, words: List[str]) -> str:
        adj = {c: set() for w in words for c in w}
        indeg = {c: 0 for c in adj}
        for w1, w2 in zip(words, words[1:]):
            if len(w1) > len(w2) and w1.startswith(w2):
                return ""
            for a, b in zip(w1, w2):
                if a != b:
                    if b not in adj[a]:
                        adj[a].add(b)
                        indeg[b] += 1
                    break
        q = deque(c for c in indeg if indeg[c] == 0)
        res = []
        while q:
            c = q.popleft()
            res.append(c)
            for nb in adj[c]:
                indeg[nb] -= 1
                if indeg[nb] == 0:
                    q.append(nb)
        return "".join(res) if len(res) == len(adj) else ""
\end{code}
\end{problem}

\begin{problem}{787}{Cheapest Flights Within K Stops}{Medium}
\summary{Cheapest price from src to dst using at most k stops (-1 if none).}
\keyidea{Bellman-Ford limited to k+1 rounds; update from a copy of last round's prices.}
\complexity{O(k * E) time, O(n) space}
\begin{code}
class Solution:
    def findCheapestPrice(self, n: int, flights: List[List[int]], src: int, dst: int, k: int) -> int:
        prices = [float("inf")] * n
        prices[src] = 0
        for _ in range(k + 1):
            tmp = prices[:]
            for u, v, w in flights:
                if prices[u] != float("inf") and prices[u] + w < tmp[v]:
                    tmp[v] = prices[u] + w
            prices = tmp
        return -1 if prices[dst] == float("inf") else prices[dst]
\end{code}
\end{problem}

\newpage
% ===== 1-D Dynamic Programming (12 problems) =====
\section{1-D Dynamic Programming \texorpdfstring{{\normalsize\color{muted}$\cdot$ 12 problems}}{- 12 problems}}
\markboth{1-D Dynamic Programming}{}

\begin{problem}{70}{Climbing Stairs}{Easy}
\summary{Ways to climb n steps taking 1 or 2 at a time.}
\keyidea{ways(n) = ways(n-1) + ways(n-2): Fibonacci with two variables.}
\complexity{O(n) time, O(1) space}
\begin{code}
class Solution:
    def climbStairs(self, n: int) -> int:
        a, b = 1, 1
        for _ in range(n - 1):
            a, b = b, a + b
        return b
\end{code}
\end{problem}

\begin{problem}{746}{Min Cost Climbing Stairs}{Easy}
\summary{Pay cost[i] to step from stair i (1 or 2 steps); start at 0 or 1; min cost to reach the top.}
\keyidea{dp[i] = cost[i] + min(dp[i+1], dp[i+2]); answer min(dp[0], dp[1]).}
\complexity{O(n) time, O(1) space}
\begin{code}
class Solution:
    def minCostClimbingStairs(self, cost: List[int]) -> int:
        a = b = 0                       # dp[i+1], dp[i+2]
        for i in range(len(cost) - 1, -1, -1):
            a, b = cost[i] + min(a, b), a
        return min(a, b)
\end{code}
\end{problem}

\begin{problem}{198}{House Robber}{Medium}
\summary{Max money from houses in a line without robbing two adjacent ones.}
\keyidea{best(i) = max(best(i-1), best(i-2) + nums[i]).}
\complexity{O(n) time, O(1) space}
\begin{code}
class Solution:
    def rob(self, nums: List[int]) -> int:
        prev2 = prev1 = 0
        for n in nums:
            prev2, prev1 = prev1, max(prev1, prev2 + n)
        return prev1
\end{code}
\end{problem}

\begin{problem}{213}{House Robber II}{Medium}
\summary{Same as House Robber but houses form a circle (first and last are adjacent).}
\keyidea{Run House Robber on nums[1:] and nums[:-1], take the max (handle n == 1).}
\complexity{O(n) time, O(1) space}
\begin{code}
class Solution:
    def rob(self, nums: List[int]) -> int:
        def line(arr):
            a = b = 0
            for n in arr:
                a, b = b, max(b, a + n)
            return b
        if len(nums) == 1:
            return nums[0]
        return max(line(nums[1:]), line(nums[:-1]))
\end{code}
\end{problem}

\begin{problem}{5}{Longest Palindromic Substring}{Medium}
\summary{Return the longest palindromic substring.}
\keyidea{Expand around every centre (odd: i,i and even: i,i+1).}
\complexity{O(n\textasciicircum{}2) time, O(1) space}
\begin{code}
class Solution:
    def longestPalindrome(self, s: str) -> str:
        res = ""
        for i in range(len(s)):
            for l, r in ((i, i), (i, i + 1)):
                while l >= 0 and r < len(s) and s[l] == s[r]:
                    l, r = l - 1, r + 1
                if r - l - 1 > len(res):
                    res = s[l + 1 : r]
        return res
\end{code}
\end{problem}

\begin{problem}{647}{Palindromic Substrings}{Medium}
\summary{Count all palindromic substrings.}
\keyidea{Expand around every centre; each successful expansion is one more palindrome.}
\complexity{O(n\textasciicircum{}2) time, O(1) space}
\begin{code}
class Solution:
    def countSubstrings(self, s: str) -> int:
        count = 0
        for i in range(len(s)):
            for l, r in ((i, i), (i, i + 1)):
                while l >= 0 and r < len(s) and s[l] == s[r]:
                    count += 1
                    l, r = l - 1, r + 1
        return count
\end{code}
\end{problem}

\begin{problem}{91}{Decode Ways}{Medium}
\summary{Digits map 1-\textgreater{}A ... 26-\textgreater{}Z. How many ways can the digit string be decoded?}
\keyidea{dp[i] = dp[i+1] (if s[i] != '0') + dp[i+2] (if s[i:i+2] in 10..26).}
\complexity{O(n) time, O(1) space}
\begin{code}
class Solution:
    def numDecodings(self, s: str) -> int:
        one, two = 1, 0                   # dp[i+1], dp[i+2]
        for i in range(len(s) - 1, -1, -1):
            cur = 0
            if s[i] != "0":
                cur = one
                if i + 1 < len(s) and 10 <= int(s[i : i + 2]) <= 26:
                    cur += two
            one, two = cur, one
        return one
\end{code}
\end{problem}

\begin{problem}{322}{Coin Change}{Medium}
\summary{Fewest coins (unlimited supply) to make amount; -1 if impossible.}
\keyidea{dp[a] = 1 + min(dp[a - c]) over coins c (unbounded knapsack).}
\complexity{O(amount * coins) time, O(amount) space}
\begin{code}
class Solution:
    def coinChange(self, coins: List[int], amount: int) -> int:
        dp = [0] + [amount + 1] * amount
        for a in range(1, amount + 1):
            for c in coins:
                if c <= a:
                    dp[a] = min(dp[a], 1 + dp[a - c])
        return dp[amount] if dp[amount] <= amount else -1
\end{code}
\end{problem}

\begin{problem}{152}{Maximum Product Subarray}{Medium}
\summary{Largest product of a contiguous subarray.}
\keyidea{Track both current max and min product (a negative can flip min into max).}
\complexity{O(n) time, O(1) space}
\begin{code}
class Solution:
    def maxProduct(self, nums: List[int]) -> int:
        res = max(nums)
        cur_max = cur_min = 1
        for n in nums:
            cand = (n, n * cur_max, n * cur_min)
            cur_max, cur_min = max(cand), min(cand)
            res = max(res, cur_max)
        return res
\end{code}
\end{problem}

\begin{problem}{139}{Word Break}{Medium}
\summary{Can s be split into a sequence of dictionary words (reuse allowed)?}
\keyidea{dp[i] = True if some word w matches s[i:] start and dp[i + len(w)] is True.}
\complexity{O(n * m * L) time, O(n) space}
\begin{code}
class Solution:
    def wordBreak(self, s: str, wordDict: List[str]) -> bool:
        dp = [False] * (len(s) + 1)
        dp[len(s)] = True
        for i in range(len(s) - 1, -1, -1):
            for w in wordDict:
                if s.startswith(w, i) and dp[i + len(w)]:
                    dp[i] = True
                    break
        return dp[0]
\end{code}
\end{problem}

\begin{problem}{300}{Longest Increasing Subsequence}{Medium}
\summary{Length of the longest strictly increasing subsequence.}
\keyidea{Patience sorting: tails[k] = smallest tail of an increasing subsequence of length k+1; binary search to place each number.}
\complexity{O(n log n) time, O(n) space}
\begin{code}
class Solution:
    def lengthOfLIS(self, nums: List[int]) -> int:
        tails = []
        for n in nums:
            i = bisect.bisect_left(tails, n)
            if i == len(tails):
                tails.append(n)
            else:
                tails[i] = n
        return len(tails)
\end{code}
\end{problem}

\begin{problem}{416}{Partition Equal Subset Sum}{Medium}
\summary{Can the array be split into two subsets with equal sums?}
\keyidea{Need a subset summing to total/2: keep a set (or bitset) of reachable sums.}
\complexity{O(n * sum) time, O(sum) space}
\begin{code}
class Solution:
    def canPartition(self, nums: List[int]) -> bool:
        total = sum(nums)
        if total % 2:
            return False
        target = total // 2
        reach = {0}
        for n in nums:
            reach |= {r + n for r in reach if r + n <= target}
            if target in reach:
                return True
        return target in reach
\end{code}
\end{problem}

\newpage
% ===== 2-D Dynamic Programming (11 problems) =====
\section{2-D Dynamic Programming \texorpdfstring{{\normalsize\color{muted}$\cdot$ 11 problems}}{- 11 problems}}
\markboth{2-D Dynamic Programming}{}

\begin{problem}{62}{Unique Paths}{Medium}
\summary{Paths from top-left to bottom-right of an m x n grid moving only right/down.}
\keyidea{paths(cell) = paths(below) + paths(right); keep one row.}
\complexity{O(m * n) time, O(n) space}
\begin{code}
class Solution:
    def uniquePaths(self, m: int, n: int) -> int:
        row = [1] * n
        for _ in range(m - 1):
            for j in range(n - 2, -1, -1):
                row[j] += row[j + 1]
        return row[0]
\end{code}
\end{problem}

\begin{problem}{1143}{Longest Common Subsequence}{Medium}
\summary{Length of the longest subsequence common to two strings.}
\keyidea{dp[i][j] = 1 + dp[i+1][j+1] if chars match else max(dp[i+1][j], dp[i][j+1]).}
\complexity{O(m * n) time and space}
\begin{code}
class Solution:
    def longestCommonSubsequence(self, text1: str, text2: str) -> int:
        m, n = len(text1), len(text2)
        dp = [[0] * (n + 1) for _ in range(m + 1)]
        for i in range(m - 1, -1, -1):
            for j in range(n - 1, -1, -1):
                if text1[i] == text2[j]:
                    dp[i][j] = 1 + dp[i + 1][j + 1]
                else:
                    dp[i][j] = max(dp[i + 1][j], dp[i][j + 1])
        return dp[0][0]
\end{code}
\end{problem}

\begin{problem}{309}{Best Time to Buy and Sell Stock with Cooldown}{Medium}
\summary{Unlimited trades, but after selling you must wait one day before buying. Max profit.}
\keyidea{State machine: hold, sold (cooldown today), rest. Transition each day.}
\complexity{O(n) time, O(1) space}
\begin{code}
class Solution:
    def maxProfit(self, prices: List[int]) -> int:
        hold, sold, rest = float("-inf"), 0, 0
        for p in prices:
            hold, sold, rest = max(hold, rest - p), hold + p, max(rest, sold)
        return max(sold, rest)
\end{code}
\end{problem}

\begin{problem}{518}{Coin Change II}{Medium}
\summary{Number of combinations (order doesn't matter) of coins that make amount.}
\keyidea{Loop coins in the outer loop so each combination is counted once: dp[a] += dp[a - c].}
\complexity{O(amount * coins) time, O(amount) space}
\begin{code}
class Solution:
    def change(self, amount: int, coins: List[int]) -> int:
        dp = [1] + [0] * amount
        for c in coins:
            for a in range(c, amount + 1):
                dp[a] += dp[a - c]
        return dp[amount]
\end{code}
\end{problem}

\begin{problem}{494}{Target Sum}{Medium}
\summary{Put + or - before each number; count expressions that equal target.}
\keyidea{Map of reachable sum -\textgreater{} number of ways, updated number by number.}
\complexity{O(n * sum) time, O(sum) space}
\begin{code}
class Solution:
    def findTargetSumWays(self, nums: List[int], target: int) -> int:
        ways = {0: 1}
        for n in nums:
            nxt = defaultdict(int)
            for s, cnt in ways.items():
                nxt[s + n] += cnt
                nxt[s - n] += cnt
            ways = nxt
        return ways.get(target, 0)
\end{code}
\end{problem}

\begin{problem}{97}{Interleaving String}{Medium}
\summary{Is s3 formed by interleaving s1 and s2 while keeping each one's order?}
\keyidea{dp[i][j] = can s1[i:] and s2[j:] form s3[i+j:].}
\complexity{O(m * n) time and space}
\begin{code}
class Solution:
    def isInterleave(self, s1: str, s2: str, s3: str) -> bool:
        m, n = len(s1), len(s2)
        if m + n != len(s3):
            return False
        dp = [[False] * (n + 1) for _ in range(m + 1)]
        dp[m][n] = True
        for i in range(m, -1, -1):
            for j in range(n, -1, -1):
                if i < m and s1[i] == s3[i + j] and dp[i + 1][j]:
                    dp[i][j] = True
                if j < n and s2[j] == s3[i + j] and dp[i][j + 1]:
                    dp[i][j] = True
        return dp[0][0]
\end{code}
\end{problem}

\begin{problem}{329}{Longest Increasing Path in a Matrix}{Hard}
\summary{Longest strictly increasing path moving in 4 directions.}
\keyidea{DFS + memo: lip(cell) = 1 + max lip(neighbour with larger value). No visited set needed (strictly increasing = no cycles).}
\complexity{O(m * n) time and space}
\begin{code}
class Solution:
    def longestIncreasingPath(self, matrix: List[List[int]]) -> int:
        R, C = len(matrix), len(matrix[0])
        @functools.lru_cache(None)
        def dfs(r, c):
            best = 1
            for nr, nc in ((r + 1, c), (r - 1, c), (r, c + 1), (r, c - 1)):
                if 0 <= nr < R and 0 <= nc < C and matrix[nr][nc] > matrix[r][c]:
                    best = max(best, 1 + dfs(nr, nc))
            return best
        return max(dfs(r, c) for r in range(R) for c in range(C))
\end{code}
\end{problem}

\begin{problem}{115}{Distinct Subsequences}{Hard}
\summary{How many distinct subsequences of s equal t?}
\keyidea{dp[i][j] = dp[i+1][j] (skip s[i]) + dp[i+1][j+1] (use it, if s[i] == t[j]).}
\complexity{O(m * n) time and space}
\begin{code}
class Solution:
    def numDistinct(self, s: str, t: str) -> int:
        m, n = len(s), len(t)
        dp = [[0] * (n + 1) for _ in range(m + 1)]
        for i in range(m + 1):
            dp[i][n] = 1
        for i in range(m - 1, -1, -1):
            for j in range(n - 1, -1, -1):
                dp[i][j] = dp[i + 1][j]
                if s[i] == t[j]:
                    dp[i][j] += dp[i + 1][j + 1]
        return dp[0][0]
\end{code}
\end{problem}

\begin{problem}{72}{Edit Distance}{Hard}
\summary{Minimum inserts, deletes and replacements to turn word1 into word2.}
\keyidea{If chars match, dp[i+1][j+1]; else 1 + min(insert dp[i][j+1], delete dp[i+1][j], replace dp[i+1][j+1]).}
\complexity{O(m * n) time and space}
\begin{code}
class Solution:
    def minDistance(self, word1: str, word2: str) -> int:
        m, n = len(word1), len(word2)
        dp = [[0] * (n + 1) for _ in range(m + 1)]
        for i in range(m + 1): dp[i][n] = m - i
        for j in range(n + 1): dp[m][j] = n - j
        for i in range(m - 1, -1, -1):
            for j in range(n - 1, -1, -1):
                if word1[i] == word2[j]:
                    dp[i][j] = dp[i + 1][j + 1]
                else:
                    dp[i][j] = 1 + min(dp[i + 1][j], dp[i][j + 1], dp[i + 1][j + 1])
        return dp[0][0]
\end{code}
\end{problem}

\begin{problem}{312}{Burst Balloons}{Hard}
\summary{Bursting balloon i earns left * i * right (current neighbours). Maximise total coins.}
\keyidea{Interval DP: choose which balloon in (l, r) is burst LAST; its neighbours are then fixed as l and r.}
\complexity{O(n\textasciicircum{}3) time, O(n\textasciicircum{}2) space}
\begin{code}
class Solution:
    def maxCoins(self, nums: List[int]) -> int:
        a = [1] + nums + [1]
        n = len(a)
        dp = [[0] * n for _ in range(n)]
        for length in range(2, n):
            for l in range(0, n - length):
                r = l + length
                for k in range(l + 1, r):
                    dp[l][r] = max(dp[l][r], a[l] * a[k] * a[r] + dp[l][k] + dp[k][r])
        return dp[0][n - 1]
\end{code}
\end{problem}

\begin{problem}{10}{Regular Expression Matching}{Hard}
\summary{Full-string match supporting '.' (any char) and '*' (zero or more of the previous char).}
\keyidea{Top-down DP on (i, j). If p[j+1] == '*': skip the pair, or use it if first chars match.}
\complexity{O(m * n) time and space}
\begin{code}
class Solution:
    def isMatch(self, s: str, p: str) -> bool:
        @functools.lru_cache(None)
        def dp(i, j):
            if j == len(p):
                return i == len(s)
            first = i < len(s) and p[j] in (s[i], ".")
            if j + 1 < len(p) and p[j + 1] == "*":
                return dp(i, j + 2) or (first and dp(i + 1, j))
            return first and dp(i + 1, j + 1)
        return dp(0, 0)
\end{code}
\end{problem}

\newpage
% ===== Greedy (8 problems) =====
\section{Greedy {\normalsize\color{muted}$\cdot$ 8 problems}}
\markboth{Greedy}{}

\begin{problem}{53}{Maximum Subarray}{Medium}
\summary{Largest sum of a contiguous subarray.}
\keyidea{Kadane: drop the running sum when it goes negative.}
\complexity{O(n) time, O(1) space}
\begin{code}
class Solution:
    def maxSubArray(self, nums: List[int]) -> int:
        best, cur = nums[0], 0
        for n in nums:
            cur = max(cur, 0) + n
            best = max(best, cur)
        return best
\end{code}
\end{problem}

\begin{problem}{55}{Jump Game}{Medium}
\summary{nums[i] is the max jump from i. Can you reach the last index?}
\keyidea{Walk backwards moving a 'goal' left whenever i + nums[i] \textgreater{}= goal.}
\complexity{O(n) time, O(1) space}
\begin{code}
class Solution:
    def canJump(self, nums: List[int]) -> bool:
        goal = len(nums) - 1
        for i in range(len(nums) - 2, -1, -1):
            if i + nums[i] >= goal:
                goal = i
        return goal == 0
\end{code}
\end{problem}

\begin{problem}{45}{Jump Game II}{Medium}
\summary{Minimum jumps to reach the last index (always reachable).}
\keyidea{BFS by windows: [l, r] is reachable in 'jumps'; next window ends at the farthest reach.}
\complexity{O(n) time, O(1) space}
\begin{code}
class Solution:
    def jump(self, nums: List[int]) -> int:
        jumps = l = r = 0
        while r < len(nums) - 1:
            farthest = max(i + nums[i] for i in range(l, r + 1))
            l, r = r + 1, farthest
            jumps += 1
        return jumps
\end{code}
\end{problem}

\begin{problem}{134}{Gas Station}{Medium}
\summary{Circular route with gas[i] and cost[i]; find the unique start index to complete the loop, or -1.}
\keyidea{If total gas \textless{} total cost: -1. Otherwise reset the start to i+1 whenever the tank goes negative.}
\complexity{O(n) time, O(1) space}
\begin{code}
class Solution:
    def canCompleteCircuit(self, gas: List[int], cost: List[int]) -> int:
        if sum(gas) < sum(cost):
            return -1
        tank = start = 0
        for i in range(len(gas)):
            tank += gas[i] - cost[i]
            if tank < 0:
                tank, start = 0, i + 1
        return start
\end{code}
\end{problem}

\begin{problem}{846}{Hand of Straights}{Medium}
\summary{Can the cards be split into groups of groupSize consecutive values?}
\keyidea{Always start a group at the smallest remaining card (min-heap / sorted keys) and consume the run.}
\complexity{O(n log n) time, O(n) space}
\begin{code}
class Solution:
    def isNStraightHand(self, hand: List[int], groupSize: int) -> bool:
        if len(hand) % groupSize:
            return False
        count = Counter(hand)
        for start in sorted(count):
            need = count[start]
            if need:
                for v in range(start, start + groupSize):
                    if count[v] < need:
                        return False
                    count[v] -= need
        return True
\end{code}
\end{problem}

\begin{problem}{1899}{Merge Triplets to Form Target Triplet}{Medium}
\summary{Merging takes element-wise max. Can some merges produce exactly the target triplet?}
\keyidea{Ignore any triplet exceeding target anywhere; from the rest, check every target position is matched by someone.}
\complexity{O(n) time, O(1) space}
\begin{code}
class Solution:
    def mergeTriplets(self, triplets: List[List[int]], target: List[int]) -> bool:
        good = set()
        for t in triplets:
            if t[0] > target[0] or t[1] > target[1] or t[2] > target[2]:
                continue
            for i in range(3):
                if t[i] == target[i]:
                    good.add(i)
        return len(good) == 3
\end{code}
\end{problem}

\begin{problem}{763}{Partition Labels}{Medium}
\summary{Split s into as many parts as possible so each letter appears in only one part; return part sizes.}
\keyidea{Record each letter's last index; extend the current part's end to the max last index seen; cut when i == end.}
\complexity{O(n) time, O(26) space}
\begin{code}
class Solution:
    def partitionLabels(self, s: str) -> List[int]:
        last = {c: i for i, c in enumerate(s)}
        res, size, end = [], 0, 0
        for i, c in enumerate(s):
            size += 1
            end = max(end, last[c])
            if i == end:
                res.append(size)
                size = 0
        return res
\end{code}
\end{problem}

\begin{problem}{678}{Valid Parenthesis String}{Medium}
\summary{'(' , ')' and '*' (which may be '(', ')' or empty). Can the string be valid?}
\keyidea{Track a range [lo, hi] of possible open counts; hi \textless{} 0 -\textgreater{} invalid; clamp lo at 0; valid if lo == 0 at end.}
\complexity{O(n) time, O(1) space}
\begin{code}
class Solution:
    def checkValidString(self, s: str) -> bool:
        lo = hi = 0
        for c in s:
            if c == "(":
                lo, hi = lo + 1, hi + 1
            elif c == ")":
                lo, hi = lo - 1, hi - 1
            else:
                lo, hi = lo - 1, hi + 1
            if hi < 0:
                return False
            lo = max(lo, 0)
        return lo == 0
\end{code}
\end{problem}

\newpage
% ===== Intervals (6 problems) =====
\section{Intervals \texorpdfstring{{\normalsize\color{muted}$\cdot$ 6 problems}}{- 6 problems}}
\markboth{Intervals}{}

\begin{problem}{57}{Insert Interval}{Medium}
\summary{Insert a new interval into a sorted, non-overlapping list and merge as needed.}
\keyidea{Three phases: copy intervals ending before it; merge all overlapping into it; copy the rest.}
\complexity{O(n) time, O(n) space}
\begin{code}
class Solution:
    def insert(self, intervals: List[List[int]], newInterval: List[int]) -> List[List[int]]:
        res, i, n = [], 0, len(intervals)
        while i < n and intervals[i][1] < newInterval[0]:
            res.append(intervals[i]); i += 1
        while i < n and intervals[i][0] <= newInterval[1]:
            newInterval = [min(newInterval[0], intervals[i][0]),
                           max(newInterval[1], intervals[i][1])]
            i += 1
        res.append(newInterval)
        res.extend(intervals[i:])
        return res
\end{code}
\end{problem}

\begin{problem}{56}{Merge Intervals}{Medium}
\summary{Merge all overlapping intervals.}
\keyidea{Sort by start; extend the last merged interval if the next one starts before it ends.}
\complexity{O(n log n) time, O(n) space}
\begin{code}
class Solution:
    def merge(self, intervals: List[List[int]]) -> List[List[int]]:
        intervals.sort(key=lambda x: x[0])
        res = [intervals[0][:]]
        for s, e in intervals[1:]:
            if s <= res[-1][1]:
                res[-1][1] = max(res[-1][1], e)
            else:
                res.append([s, e])
        return res
\end{code}
\end{problem}

\begin{problem}{435}{Non-overlapping Intervals}{Medium}
\summary{Minimum intervals to remove so the rest don't overlap.}
\keyidea{Sort by start; on overlap, remove the one with the larger end (keep the earlier-ending one).}
\complexity{O(n log n) time, O(1) space}
\begin{code}
class Solution:
    def eraseOverlapIntervals(self, intervals: List[List[int]]) -> int:
        intervals.sort()
        removed, prev_end = 0, intervals[0][1]
        for s, e in intervals[1:]:
            if s >= prev_end:
                prev_end = e
            else:
                removed += 1
                prev_end = min(prev_end, e)
        return removed
\end{code}
\end{problem}

\begin{problem}{252}{Meeting Rooms}{Easy}
\summary{Can one person attend all meetings (no overlaps)?}
\keyidea{Sort by start; any meeting starting before the previous ends is a conflict.}
\complexity{O(n log n) time, O(1) space}
\begin{code}
class Solution:
    def canAttendMeetings(self, intervals: List[List[int]]) -> bool:
        intervals.sort()
        for i in range(1, len(intervals)):
            if intervals[i][0] < intervals[i - 1][1]:
                return False
        return True
\end{code}
\end{problem}

\begin{problem}{253}{Meeting Rooms II}{Medium}
\summary{Minimum number of rooms needed for all meetings.}
\keyidea{Sort starts and ends separately; sweep: a start before the earliest end needs a new room.}
\complexity{O(n log n) time, O(n) space}
\begin{code}
class Solution:
    def minMeetingRooms(self, intervals: List[List[int]]) -> int:
        starts = sorted(i[0] for i in intervals)
        ends = sorted(i[1] for i in intervals)
        rooms = best = e = 0
        for s in starts:
            while e < len(ends) and ends[e] <= s:
                rooms -= 1
                e += 1
            rooms += 1
            best = max(best, rooms)
        return best
\end{code}
\end{problem}

\begin{problem}{1851}{Minimum Interval to Include Each Query}{Hard}
\summary{For each query point, return the size of the smallest interval containing it (-1 if none).}
\keyidea{Process queries sorted; push intervals whose start \textless{}= q into a min-heap by size; pop those ending before q.}
\complexity{O((n + q) log n) time, O(n + q) space}
\begin{code}
class Solution:
    def minInterval(self, intervals: List[List[int]], queries: List[int]) -> List[int]:
        intervals.sort()
        h, ans, i = [], {}, 0
        for q in sorted(queries):
            while i < len(intervals) and intervals[i][0] <= q:
                l, r = intervals[i]
                heapq.heappush(h, (r - l + 1, r))
                i += 1
            while h and h[0][1] < q:
                heapq.heappop(h)
            ans[q] = h[0][0] if h else -1
        return [ans[q] for q in queries]
\end{code}
\end{problem}

\newpage
% ===== Math & Geometry (8 problems) =====
\section{Math \& Geometry {\normalsize\color{muted}$\cdot$ 8 problems}}
\markboth{Math \& Geometry}{}

\begin{problem}{48}{Rotate Image}{Medium}
\summary{Rotate an n x n matrix 90 degrees clockwise in place.}
\keyidea{Transpose, then reverse each row.}
\complexity{O(n\textasciicircum{}2) time, O(1) space}
\begin{code}
class Solution:
    def rotate(self, matrix: List[List[int]]) -> None:
        n = len(matrix)
        for i in range(n):
            for j in range(i + 1, n):
                matrix[i][j], matrix[j][i] = matrix[j][i], matrix[i][j]
        for row in matrix:
            row.reverse()
\end{code}
\end{problem}

\begin{problem}{54}{Spiral Matrix}{Medium}
\summary{Return all elements in clockwise spiral order.}
\keyidea{Shrink four boundaries (top, bottom, left, right) after each side is read.}
\complexity{O(m * n) time, O(1) extra space}
\begin{code}
class Solution:
    def spiralOrder(self, matrix: List[List[int]]) -> List[int]:
        res = []
        top, bottom, left, right = 0, len(matrix) - 1, 0, len(matrix[0]) - 1
        while top <= bottom and left <= right:
            for c in range(left, right + 1): res.append(matrix[top][c])
            top += 1
            for r in range(top, bottom + 1): res.append(matrix[r][right])
            right -= 1
            if top <= bottom:
                for c in range(right, left - 1, -1): res.append(matrix[bottom][c])
                bottom -= 1
            if left <= right:
                for r in range(bottom, top - 1, -1): res.append(matrix[r][left])
                left += 1
        return res
\end{code}
\end{problem}

\begin{problem}{73}{Set Matrix Zeroes}{Medium}
\summary{If a cell is 0, set its entire row and column to 0, in place.}
\keyidea{Use the first row/column as marker arrays; keep a separate flag for row 0.}
\complexity{O(m * n) time, O(1) space}
\begin{code}
class Solution:
    def setZeroes(self, matrix: List[List[int]]) -> None:
        R, C = len(matrix), len(matrix[0])
        row0_zero = any(matrix[0][c] == 0 for c in range(C))
        for r in range(1, R):
            for c in range(C):
                if matrix[r][c] == 0:
                    matrix[0][c] = 0
                    matrix[r][0] = 0
        for r in range(1, R):
            for c in range(1, C):
                if matrix[0][c] == 0 or matrix[r][0] == 0:
                    matrix[r][c] = 0
        if matrix[0][0] == 0:
            for r in range(R):
                matrix[r][0] = 0
        if row0_zero:
            for c in range(C):
                matrix[0][c] = 0
\end{code}
\end{problem}

\begin{problem}{202}{Happy Number}{Easy}
\summary{Repeatedly replace n with the sum of squares of its digits. Does it reach 1?}
\keyidea{Either reaches 1 or loops: detect the loop with a set (or Floyd).}
\complexity{O(log n) per step, O(log n) space}
\begin{code}
class Solution:
    def isHappy(self, n: int) -> bool:
        seen = set()
        while n != 1 and n not in seen:
            seen.add(n)
            n = sum(int(d) ** 2 for d in str(n))
        return n == 1
\end{code}
\end{problem}

\begin{problem}{66}{Plus One}{Easy}
\summary{Add one to a number stored as a list of digits.}
\keyidea{From the end: 9 becomes 0 and carry; otherwise increment and return. All 9s -\textgreater{} prepend 1.}
\complexity{O(n) time, O(1) space}
\begin{code}
class Solution:
    def plusOne(self, digits: List[int]) -> List[int]:
        for i in range(len(digits) - 1, -1, -1):
            if digits[i] < 9:
                digits[i] += 1
                return digits
            digits[i] = 0
        return [1] + digits
\end{code}
\end{problem}

\begin{problem}{50}{Pow(x, n)}{Medium}
\summary{Compute x raised to integer n (n may be negative).}
\keyidea{Fast exponentiation: square x and halve n; multiply in x when the low bit is 1.}
\complexity{O(log n) time, O(1) space}
\begin{code}
class Solution:
    def myPow(self, x: float, n: int) -> float:
        if n < 0:
            x, n = 1 / x, -n
        res = 1.0
        while n:
            if n & 1:
                res *= x
            x *= x
            n >>= 1
        return res
\end{code}
\end{problem}

\begin{problem}{43}{Multiply Strings}{Medium}
\summary{Multiply two non-negative numbers given as strings, without converting to int.}
\keyidea{Grade-school multiplication: digit i times digit j lands at position i + j (on reversed strings).}
\complexity{O(m * n) time, O(m + n) space}
\begin{code}
class Solution:
    def multiply(self, num1: str, num2: str) -> str:
        if num1 == "0" or num2 == "0":
            return "0"
        res = [0] * (len(num1) + len(num2))
        a, b = num1[::-1], num2[::-1]
        for i in range(len(a)):
            for j in range(len(b)):
                res[i + j] += int(a[i]) * int(b[j])
                res[i + j + 1] += res[i + j] // 10
                res[i + j] %= 10
        while len(res) > 1 and res[-1] == 0:
            res.pop()
        return "".join(map(str, res[::-1]))
\end{code}
\end{problem}

\begin{problem}{2013}{Detect Squares}{Medium}
\summary{add(point) to a multiset; count(point) returns how many axis-aligned squares it completes with stored points.}
\keyidea{Treat each stored point as the diagonal corner; if it forms a square, multiply counts of the other two corners.}
\complexity{add O(1), count O(n) time; O(n) space}
\begin{code}
class DetectSquares:
    def __init__(self):
        self.pts = Counter()

    def add(self, point: List[int]) -> None:
        self.pts[tuple(point)] += 1

    def count(self, point: List[int]) -> int:
        px, py = point
        res = 0
        for (x, y), c in list(self.pts.items()):
            if abs(px - x) != abs(py - y) or x == px or y == py:
                continue
            res += c * self.pts[(x, py)] * self.pts[(px, y)]
        return res
\end{code}
\end{problem}

\newpage
% ===== Bit Manipulation (7 problems) =====
\section{Bit Manipulation \texorpdfstring{{\normalsize\color{muted}$\cdot$ 7 problems}}{- 7 problems}}
\markboth{Bit Manipulation}{}

\begin{problem}{136}{Single Number}{Easy}
\summary{Every value appears twice except one; find it in O(1) space.}
\keyidea{XOR everything: pairs cancel (a \textasciicircum{} a = 0).}
\complexity{O(n) time, O(1) space}
\begin{code}
class Solution:
    def singleNumber(self, nums: List[int]) -> int:
        res = 0
        for n in nums:
            res ^= n
        return res
\end{code}
\end{problem}

\begin{problem}{191}{Number of 1 Bits}{Easy}
\summary{Count the set bits of an unsigned integer.}
\keyidea{n \&= n - 1 clears the lowest set bit; count iterations.}
\complexity{O(number of set bits) time, O(1) space}
\begin{code}
class Solution:
    def hammingWeight(self, n: int) -> int:
        count = 0
        while n:
            n &= n - 1
            count += 1
        return count
\end{code}
\end{problem}

\begin{problem}{338}{Counting Bits}{Easy}
\summary{For every i in 0..n, return the number of 1 bits in i.}
\keyidea{dp[i] = dp[i \textgreater{}\textgreater{} 1] + (i \& 1).}
\complexity{O(n) time, O(n) space}
\begin{code}
class Solution:
    def countBits(self, n: int) -> List[int]:
        dp = [0] * (n + 1)
        for i in range(1, n + 1):
            dp[i] = dp[i >> 1] + (i & 1)
        return dp
\end{code}
\end{problem}

\begin{problem}{190}{Reverse Bits}{Easy}
\summary{Reverse the bits of a 32-bit unsigned integer.}
\keyidea{Take bit i from n and place it at position 31 - i.}
\complexity{O(32) time, O(1) space}
\begin{code}
class Solution:
    def reverseBits(self, n: int) -> int:
        res = 0
        for i in range(32):
            bit = (n >> i) & 1
            res |= bit << (31 - i)
        return res
\end{code}
\end{problem}

\begin{problem}{268}{Missing Number}{Easy}
\summary{Array holds n distinct numbers from 0..n; find the missing one.}
\keyidea{XOR all indices 0..n with all values (or n(n+1)/2 - sum).}
\complexity{O(n) time, O(1) space}
\begin{code}
class Solution:
    def missingNumber(self, nums: List[int]) -> int:
        res = len(nums)
        for i, n in enumerate(nums):
            res ^= i ^ n
        return res
\end{code}
\end{problem}

\begin{problem}{371}{Sum of Two Integers}{Medium}
\summary{Add two integers without using + or -.}
\keyidea{sum = a \textasciicircum{} b, carry = (a \& b) \textless{}\textless{} 1; repeat. In Python, mask to 32 bits and fix the sign at the end.}
\complexity{O(32) time, O(1) space}
\begin{code}
class Solution:
    def getSum(self, a: int, b: int) -> int:
        mask = 0xFFFFFFFF
        while b & mask:
            a, b = a ^ b, (a & b) << 1
        a &= mask
        return a if a <= 0x7FFFFFFF else ~(a ^ mask)
\end{code}
\end{problem}

\begin{problem}{7}{Reverse Integer}{Medium}
\summary{Reverse the digits of a 32-bit signed integer; return 0 if the result overflows.}
\keyidea{Pop digits with \% and //, building the result; check the 32-bit bounds.}
\complexity{O(log x) time, O(1) space}
\begin{code}
class Solution:
    def reverse(self, x: int) -> int:
        sign = -1 if x < 0 else 1
        x, res = abs(x), 0
        while x:
            res = res * 10 + x % 10
            x //= 10
        res *= sign
        return res if -2**31 <= res <= 2**31 - 1 else 0
\end{code}
\end{problem}


\end{document}
